\documentclass[a4paper, 12pt]{article}
\usepackage[utf8]{inputenc}
\usepackage[T1]{fontenc}
\usepackage{amsmath,amssymb,amsthm}
% Margins per HSE technical requirements: left 35mm, right >=10mm, top/bottom >=20mm
\usepackage[a4paper,left=35mm,right=10mm,top=20mm,bottom=20mm]{geometry}
\usepackage[english]{babel}
\usepackage{tabularx}
\usepackage{adjustbox}
\usepackage{graphicx}
\usepackage{subcaption}
\usepackage{booktabs}
\usepackage[most]{tcolorbox}
\usepackage{xcolor}
\usepackage{setspace}
\usepackage{fancyhdr}
\setlength{\headheight}{15pt}
\usepackage{etoolbox}
\usepackage{hyperref}
\usepackage{fontawesome5}

% --- APA STYLE ---
\usepackage{csquotes}
\usepackage[backend=biber, style=apa]{biblatex}
\addbibresource{bibliography.bib}
\setlength{\bibhang}{2em}

% Line spacing 1.5 for main text and reference list
\onehalfspacing
% Tables, figures and footnotes single-spaced (setspace defaults float content
% to single spacing; we explicitly enforce it for floats and footnotes below).
\AtBeginEnvironment{table}{\singlespacing}
\AtBeginEnvironment{figure}{\singlespacing}
\AtBeginEnvironment{tabular}{\singlespacing}

% Footnotes: 10pt, single-spaced
\let\origfootnote\footnote
\renewcommand{\footnote}[1]{\origfootnote{\setstretch{1.0}\fontsize{10pt}{12pt}\selectfont #1}}

% Paragraph indent 1.25 cm (justified by default; no \RaggedRight)
\setlength{\parindent}{1.25cm}

\usepackage{titlesec}

% No section numbering (manual hierarchy retained for centred display)
\setcounter{secnumdepth}{0}

% Each new \section starts on a new page (HSE rule)
\newcommand{\sectionbreak}{\clearpage}

% Level 1 (\section) — centred, bold
\titleformat{\section}
  {\normalfont\Large\bfseries\centering}{\thesection}{1em}{}

% Level 2 (\subsection)
\titleformat{\subsection}
  {\normalfont\large\bfseries}{\thesubsection}{1em}{}

% Level 3 (\subsubsection) — run-in
\titleformat{\subsubsection}[runin]
  {\normalfont\normalsize\bfseries}{\thesubsubsection}{1em}{}[.]

\titlespacing*{\subsubsection}
  {1.25cm}{1ex plus .2ex minus .2ex}{1em}

% Page numbering: bottom-centred, hidden on the title page
\pagestyle{plain}

% Appendix header in upper-right corner (set per appendix via \appendixhead)
\fancypagestyle{appendix}{%
  \fancyhf{}%
  \fancyhead[R]{\appendixname\ \theappendixnum}%
  \fancyfoot[C]{\thepage}%
  \renewcommand{\headrulewidth}{0pt}%
}
\newcounter{appendixnum}
\providecommand{\appendixname}{Appendix}
\newcommand{\appendixhead}[1]{%
  \clearpage
  \refstepcounter{appendixnum}%
  \thispagestyle{appendix}%
  \pagestyle{appendix}%
  \phantomsection
  \section*{\appendixname\ \theappendixnum: #1}%
  \addcontentsline{toc}{section}{\appendixname\ \theappendixnum: #1}%
}
% Clean, Academic Return Button
\newcommand{\goback}[1]{%
    \hfill \hyperlink{#1}{%  \hfill pushes it to the right side of the table
        \begin{tcolorbox}[
            on line, 
            boxrule=0.5pt, 
            colback=white,      % Professional white background
            colframe=gray!80,   % Subtle gray border
            arc=0pt,            % Sharp corners look more academic
            boxsep=0pt,
            left=4pt,right=4pt,top=2pt,bottom=2pt
        ]
        \scriptsize\sffamily\color{black} $\uparrow$ BACK TO TEXT
        \end{tcolorbox}%
    }%
}

\begin{document}

%%
%% Title page
%%
\begin{center}

Federal State Autonomous Educational Institution for Higher Education\\
«NATIONAL RESEARCH UNIVERSITY\\
«HIGHER SCHOOL OF ECONOMICS»\\[2ex]
Faculty «St.Petersburg School of Economics and Management»\par

\vspace{1.5cm}

{\large Ivan D. Zamyatin\par}

\vspace{1cm}

{\Large\bfseries
INNOVATION AND PRODUCTIVITY: \\ 
AN EMPIRICAL ANALYSIS OF RUSSIAN FIRMS
\par}

\vspace{0.8cm}

BACHELOR THESIS\par

\vspace{0.8cm}

% Оборачиваем в minipage, чтобы выровнять текст по левому краю, 
% но сам блок оставить по центру страницы
\makebox[\textwidth][c]{%
    \begin{minipage}{0.55\textwidth}
        Field of study: 38.03.01 Economics\\[1ex]
        Specialization: Quantitative Economics\\[1ex]
        Format: Article
    \end{minipage}
}\par

\vspace{1cm}

EDUCATIONAL PROGRAMME\\[1ex]
«International Programme in Business and Economics»\par

\vspace{0.5cm}

Student’s group number №2201\par


\vfill

% Блок руководителя с линией для подписи
\hfill
\parbox[t]{0.4\textwidth}{%
    \underline{Adviser}\\[1ex]
    Senior Lecturer\\[1ex]
    Dmitry S. Tereshchenko\\[1ex]
}%

\vfill\vfill

Saint Petersburg\\[1ex]
2026
\end{center}
\thispagestyle{empty}
\clearpage

%%
%% Table of Contents
%%
{\singlespacing
\tableofcontents
}
\clearpage

%%
%% ===========================================================================
%%

\section{Declaration of the Use of Generative Models}

\paragraph{Generative models used.}
\begin{itemize}
    \item \textbf{Anthropic Claude} (Claude Sonnet and Claude Opus, accessed primarily through the Claude Code) - \url{https://www.anthropic.com}; \url{https://claude.com/claude-code}.
    \item \textbf{Google Gemini} (web interface) - \url{https://gemini.google.com}.
\end{itemize}

\paragraph{Application by part of the thesis.}
Following the three usage categories explicitly enumerated as (i)~fully generated text, (ii)~generated text subject to editing by the author, (iii)~text written by the author with the use of generated materials - the application of AI in this thesis is described section by section below.

\begin{enumerate}
    \item \textbf{Abstract.} Drafted by the author and edited with Anthropic Claude for length, tone, and integration of the regenerated headline figures. \emph{Goal: produce a concise, regulation-compliant abstract that reflects the final empirical estimates. Means of application: category (iii).} \emph{Effectiveness: high} for routine length and consistency edits.

    \item \textbf{Introduction.} Written by the author; the text was subsequently edited with Google Gemini for grammar, register, and clarity. The substantive narrative (research question, motivation, contribution statement) is the author's own. \emph{Goal: language and stylistic refinement of an author-drafted introduction. Means of application: category (iii).} \emph{Effectiveness: high} on routine grammar and register editing; no substantive content was generated by the model.

    \item \textbf{Literature Review.} Written by the author and edited with Google Gemini for English-language polishing; Gemini was additionally used as an auxiliary search aid for identifying candidate references, every one of which was independently verified by the author against the original publication before citation. \emph{Goal: language polish on an author-written narrative and triage of candidate references. Means of application: category (iii).} \emph{Effectiveness: moderate-to-high} for editorial polish and search-prompt triage; \emph{low} (and intentionally so) on substantive synthesis, which remained the author's responsibility.

    \item \textbf{Methods.} Drafted by the author and refined with Anthropic Claude. AI assistance covered the rewriting of the structural-control-function paragraphs, the introduction of the Levinsohn-Petrin specification with R\&D as a state variable, the new \emph{identification logic of the hierarchy} paragraph in the Estimation Strategy subsection and the consistency between the prose and the equations. \emph{Goal: structural reorganization, equation-prose alignment. Means of application: category (iii).} \emph{Effectiveness: high} for multi-paragraph restructuring and notation propagation across equations.

    \item \textbf{Data and Variables.} Drafted by the author and refined with Anthropic Claude. AI assistance covered the polishing of the Descriptive Analysis subsection in line-final form, the addition of references to the corresponding figures and the synchronization of all names of variables. \emph{Goal: present descriptive evidence for the final dataset and demonstrate that the variables are suitable for production-function estimation. Means of application: category (iii).} \emph{Effectiveness: high} for table-text consistency and cross-section number propagation.

    \item \textbf{Results.} 
    Drafted by the author and refined using Anthropic Claude. AI assistance was employed to synthesize complex empirical findings and ensure precise alignment between statistical outputs, including bootstrap estimates, and their qualitative interpretation. The AI optimized the document’s structure to prioritize core findings and streamline the narrative by moving technical details to the appendices.
    \emph{Goal: To achieve a clear and structured presentation of empirical estimates while minimizing technical redundancy. Means of application: category (iii). Effectiveness: high for cross-paragraph logical consistency and the propagation of numerical values.}

    \item \textbf{Discussion.} 
    Drafted by the author and refined using Anthropic Claude. The AI supported the reorganization of the Discussion to focus on the theoretical implications of the results, facilitating the integration of findings with established economic frameworks international benchmarks. Furthermore, it aided in synthesizing multi-dimensional context blocks, including the specifics of Russian microdata, the impact of the sanctions period, and institutional implications for market efficiency.
    \emph{Goal: To align the section with academic standards for scientific discussion, prioritizing the interpretation of results through theoretical lenses over methodological recapitulation. Means of application: category (iii). Effectiveness: high for structural reorganization and the thematic integration of diverse analytical perspectives.}

    \item \textbf{Conclusion.} 
    Drafted by the author and edited using Anthropic Claude to ensure full synchronization with the final empirical estimates. AI assistance was utilized to verify the consistency, ensuring that the final synthesis accurately reflects the refined statistical evidence. The AI helped in streamlining the concluding remarks to maintain focus on the validated results while ensuring logical alignment with the research questions posed in the introduction.
    \emph{Goal: To maintain numerical and interpretive consistency across the manuscript and ensure that thesis-level claims are supported by an empirical data. Means of application: category (iii). Effectiveness: high for selective verification of quantitative claims and synthesis of concluding remarks.}

    \item \textbf{Appendices.} 
    Drafted by the author and refined using Anthropic Claude. AI assistance focused on the structural systematization of supplementary materials, including the categorization of data cleaning procedures, detailed sectoral results, and sensitivity analyses. The AI aided in optimizing the technical layout of the appendices, improving cross-referencing, ensuring the correct resolution of internal links, and reformatting tables and multi-panel figures for better readability. 
    \emph{Goal: To provide a logically organized and technically sound supplementary record of data processing and robustness checks. Means of application: category (iii) for prose and structural organization; category (ii) for the refinement of tabular data. Effectiveness: high for cross-reference integrity and overall navigational consistency.}

    \item \textbf{Python codebase.} The original code architecture and the model-implementation logic were authored by the author; Anthropic Claude was used for bug identification, refactoring of model-class helpers, and end-to-end pipeline troubleshooting. AI-proposed edits were inspected, executed, and accepted or rejected by the author through full-pipeline re-runs. \emph{Goal: technical-debugging and refactoring assistance on author-written code. Means of application: category (ii), generated text/edits subject to editing by the author, with verification through end-to-end pipeline execution.} \emph{Effectiveness: high} for unblocking specific bugs and producing immediately runnable patches; substantive modelling decisions remained the author's.
\end{enumerate}

No part of the manuscript falls under category (i): no fully generated text was incorporated into the thesis without author review and editing.

\paragraph{Compliance with the remaining principles of the HSE Declaration.}
\begin{itemize}
    \item \emph{Academic integrity.} All AI-generated artefacts (code patches, prose suggestions, translations, candidate references) were reviewed and, where necessary, corrected by the author. Numerical results were independently reproduced through end-to-end execution of the codebase; theoretical claims and citations were independently checked against original sources.
    \item \emph{Confidentiality and intellectual property.} No personal data, no confidential information beyond what is publicly available in the Russian Financial Statements Database \parencite{bondarkov2025russian}, and no third-party intellectual property were submitted to the AI systems beyond what was strictly necessary for the technical task at hand.
    \item \emph{Critical evaluation.} AI outputs were treated throughout as proposals to be evaluated, not as authoritative answers; suggestions found to be inaccurate, methodologically unsuitable, or stylistically inconsistent with the manuscript were revised or rejected by the author.
    \item \emph{Priority of human reasoning.} The analytical and interpretive core of the thesis, the formulation of the research question, the choice of econometric strategy, the structural-economic interpretation of the empirical results, and the discussion of policy implications, is the author's own work and was not delegated to generative models.
\end{itemize}


\section{Abstract}

This study investigates the causal relationship between research and development (R\&D) investment and firm-level total factor productivity (TFP) in the Russian industrial sector. Drawing on a panel of $691{,}644$ firm-year observations from the Russian Financial Statements Database for the period 2011-2023, restricted to the eight production-oriented OKVED sections, the analysis employs a hierarchical estimation strategy that progresses from reduced-form specifications through structural control-function estimators to the dynamic endogenous innovation model of \textcite{doraszelski2013r}. The results demonstrate that R\&D investment exerts a positive and statistically significant effect on productivity, although the estimated elasticity attenuates substantially once unobserved firm heterogeneity is controlled for. The structural dynamic model reveals that R\&D fundamentally reshapes the probability distribution of future productivity states, with marginal returns that are strongly decreasing in the firm's current efficiency level. Innovation investments are most productive for firms operating below the technological frontier, generating a catch-up dynamic that decays monotonically towards the frontier (from $+0.008$ at the bottom decile to $+0.001$ at the top decile of within-sample $\omega$, while remaining positive throughout). These findings challenge the assumption of uniform returns to R\&D subsidies and underscore the need for heterogeneity-aware innovation policies. In the context of a sanctioned and resource-constrained economy, the evidence suggests that targeted support for low-to-medium productivity firms yields the highest marginal efficiency gains from public R\&D incentives.


\section{Introduction}
Productivity growth constitutes a fundamental determinant of long-run economic prosperity, and understanding its drivers remains a central preoccupation of modern economics. The seminal contribution of \textcite{solow1957technical}, whose aggregate production function framework decomposed output growth into contributions from capital, labour, and a residual, total factor productivity (TFP), interpreted as technological change, established a foundational paradigm for empirical productivity analysis. However, within this framework the sources of technological progress remained exogenous and theoretically unexplained, leaving a critical gap in understanding the mechanisms through which economies sustain long-term growth.

The endogenous growth theory advanced by \textcite{romer1990endogenous} sought to address this fundamental limitation by modelling knowledge accumulation and R\&D investments as deliberate engines of growth. Notably, the primary novelty of this theoretical shift lies in the recognition that innovations arise from intentional investments in research, yielding increasing returns through knowledge spillovers. These insights bear direct relevance for estimating the contribution of R\&D and innovation to productivity not only at the macroeconomic level but also at the level of the individual firm, thereby establishing the micro-foundations for the empirical analysis pursued in this study.

Subsequently, the conceptual framework for modern data-driven research is built upon the shift from viewing technology as an external force to recognizing it as an internal process driven by individual firms. Indeed, the 2025 Nobel Prize in Economic Sciences, awarded to \textcite{mokyr2011gifts,aghion1990model} for their foundational work on innovation-driven growth, reinforces this paradigm. By synthesizing Schumpeterian 'creative destruction' with institutional history and the dynamics of competition, they demonstrate that technological progress is the fundamental driver of sustained economic evolution.

Significantly, the application of this paradigm to the Russian economy reveals particularly acute challenges that motivate the present study. Russia’s economic trajectory has been characterised by persistent resource dependency, modest productivity growth, and structural imbalances that have constrained the diversification of its industrial base. Despite possessing a immense scientific and engineering workforce, the economy has historically exhibited weak translation of research inputs into measurable productivity gains at the firm level. In particular, the post-2022 geopolitical landscape, defined by comprehensive sanctions and restricted access to global technology markets, has further intensified the imperative to understand and stimulate domestic innovation as a foundation for economic sovereignty. In this context, quantifying the returns to firm-level R\&D investment is not merely an academic exercise but a matter of direct policy relevance, as policymakers seek evidence-based levers to advance technological upgrading and industrial competitiveness.

From this problem statement, two interrelated research questions emerge that guide the empirical investigation:
\begin{enumerate}
\item How do R\&D investments affect total factor productivity in Russian industrial firms?
\item Does the productivity impact of R\&D differ across firms with varying levels of technological intensity?
\end{enumerate}

On the basis of the theoretical and empirical literature, the following hypotheses are advanced. First, R\&D investments are hypothesised to exert a positive and statistically significant effect on firm-level total factor productivity in Russian industrial firms (\textbf{H1}). This expectation is grounded in the endogenous growth framework of \textcite{romer1990endogenous} and the knowledge capital model of \textcite{griliches1979issues}, which posit that deliberate investments in research generate knowledge capital that augments productive efficiency. Second, the productivity effect of R\&D is hypothesised to be heterogeneous across technological intensity classes, with firms in high-technology sectors exhibiting larger returns to R\&D investment (\textbf{H2}). This prediction follows from the theoretical proposition that firms operating closer to the technological frontier possess greater absorptive capacity and face stronger competitive incentives to innovate \parencite{aghion1990model}.

This study examines Russian industrial firms to determine the statistical link between R\&D inputs and their resulting productivity levels. Employing panel data within an augmented production function framework, the analysis examines heterogeneous effects across firms and assesses relevant channels in the current sanctioned environment. By grounding its investigation in the Nobel-recognised growth paradigm, the research contributes empirical evidence from a significant transitional economy, offering insights directly pertinent to the formulation of technology and industrial policy aimed at fostering diversification and sustainable competitiveness.

Furthermore, the empirical identification of the link between R\&D and economic performance has long been a cornerstone of industrial organization literature, yet it remains fraught with methodological challenges. In the context of Russian transition and post-transition dynamics, the challenge is twofold: one must account for the discrete technological leaps characteristic of heavy industry modernization, while simultaneously modeling the long-term, cumulative nature of knowledge capital. At the core of this research lies the resolution of the simultaneity problem, the tendency of firms to adjust their input levels and R\&D commitments in response to productivity realizations that are known to the firm but unobserved by the researcher. Because such unobserved productivity shocks vary systematically across firms, standard regression techniques yield biased estimates of the returns to innovation. To obtain consistent estimates, a structural control-function approach is employed, influencing firm-level variation in intermediate input consumption to proxy for these heterogeneous, unobserved productivity differentials. This strategy allows the true elasticity of R\&D investments to be identified separately from confounding pro-cyclical investment effects.

Furthermore, the analysis is advanced by endogenizing the productivity process itself. Departing from the standard assumption that productivity evolves according to an exogenous Markovian drift, the framework of \textcite{doraszelski2013r} is implemented, wherein the trajectory of a firm’s TFP is actively shaped through R\&D investment. This integrated approach allows the quantification of how current innovation expenditures fundamentally reshape the probability distribution of future efficiency. The empirical analysis draws upon the Russian Financial Statements Database (RFSD) \parencite{bondarkov2025russian}, which provides a unique longitudinal perspective on the universe of Russian enterprises from 2011 to 2023, encompassing $691{,}644$ firm-year observations once the panel is restricted to the eight production-oriented OKVED sections retained for the cross-sectoral analysis. The principal findings indicate that R\&D investment exerts a positive and significant effect on firm productivity, that the marginal returns to innovation are strongly decreasing in the firm’s current efficiency level, and that the choice of estimation methodology has first-order consequences for the recovered TFP dynamics.

The remainder of this thesis is structured as follows. Section~2 reviews the principal methodological approaches to estimating the R\&D-productivity relationship. Section~3 describes the data, variable construction, and the econometric framework. Section~4 presents the empirical results. Section~5 discusses the findings in light of the existing literature. Section~6 concludes with a summary of contributions and directions for future research.

\section{Literature Review}
\label{sec:literature}
This section is dedicated to the main methodological approaches for estimating the impact of Research and Development (R\&D) on productivity. It focuses on the evolution from direct production function models, which treat knowledge as a capital stock, to more complex multi-stage frameworks that explicitly model the innovation process from investment to economic output. The discussion highlights how these methods address core econometric challenges such as endogeneity and selection bias, and summarizes key findings from their application in empirical research.

Comprehensive overviews of this field, as provided in seminal works such as \parencite{griliches1979issues,crepon1998research}, highlight the persistent challenges in addressing measurement errors, functional form specifications, and endogeneity biases. A critical gap concerns impact of spending in R\&D and its endogenous influence on productivity function. An overview of the literature reveals two main knowledge capital models: the Griliches knowledge capital model \parencite{griliches1979issues} and the Crépon-Duguet-Mairesse (CDM)\parencite{crepon1998research} model. 

The \textcite{griliches1979issues} knowledge capital model is considered to have laid the groundwork for incorporating R\&D as a distinct input in production functions. In this framework, output is modeled as a function of physical capital, labor, and accumulated knowledge capital derived from R\&D investments, typically within a Cobb-Douglas specification augmented by an exponential time trend to capture technological progress. The model is known to emphasize the estimation of elasticities or rates of return on knowledge capital, assuming separability and perfect competition. Research has shown that this approach facilitates the analysis of both level and temporal dimensions of productivity, though it is reported to be sensitive to issues such as double-counting of R\&D in physical capital and the lack of controls for spillovers or complementarities. However, the \textcite{griliches1979issues} model, whose simplicity and interpretability make it a standard in empirical studies, provides a straightforward estimate of R\&D returns \parencite{andersen2019knowledge}. For instance, extensions of the \textcite{griliches1979issues} model, as carried out by \textcite{doraszelski2013r}, have incorporated structural dynamic elements where productivity is treated as an unobserved state variable evolving endogenously through R\&D, with uncertainties and non-linearities integrated via moments-based identification. 

Notwithstanding these improvements, the model likely fails to address endogeneity in the absence of external instruments or proxies. Consequently, employing Ordinary Least Squares (OLS), while computationally simpler, yields biased coefficients, particularly because of the simultaneity between input choices and unobserved productivity shocks.

In contrast, the CDM model, developed by \textcite{crepon1998research}, is reported to extend the \textcite{griliches1979issues} framework by explicitly addressing selection into innovation and endogeneity through a sequential system of equations. This model is structured in stages: first, the decision and intensity of R\&D are modeled; second, innovation outputs (such as innovative sales or patents) are linked to R\&D inputs; and third, productivity is estimated via a production function incorporating these innovation measures. Much research has been conducted using asymptotic least squares or full information maximum likelihood for estimation, with corrections like the inverse Mills ratio applied to account for selection biases. The \textcite{crepon1998research} approach is considered to have been particularly effective for datasets with mismeasured innovation variables, as evidenced in applications by \textcite{therrien2009innovation} and \textcite{damijan2010innovation}, where firm-level panels from Canada and Slovenia, respectively, were analyzed to assess productivity growth sensitivities to sectoral heterogeneity and estimation methods. Nevertheless, ambiguities arise in the model's flexibility, which can lead to misspecification if theoretical underpinnings are not robustly integrated, such as in distinguishing between product and process innovations or incorporating externalities.

Both \textcite{griliches1979issues} and \textcite{crepon1998research} ground their models in endogenous growth theory, linking R\&D to productivity. However, \textcite{griliches1979issues} uses panel data to estimate knowledge capital in a single equation, which often overlooks R\&D selection bias. Conversely, the \textcite{crepon1998research} multi-stage CDM framework uses cross-sectional data to explicitly model this selection, improving causal inference. While the Griliches approach excels at benchmarking R\&D returns and capturing non-linearities, the CDM model provides more granular insights into innovation types and policy levers.

Following the methodological preliminaries on modeling the relationship between R\&D and productivity, we now turn to its empirical manifestations across various economies, with particular attention to emerging markets.

The overarching evidence from a substantial body of firm-, industry-, and country-level literature confirms a positive association between R\&D investment and productivity. In developed economies, panel studies from the UK and OECD countries consistently report that R\&D significantly raises total factor productivity (TFP), with estimated elasticities typically ranging between 0.07 and 0.15.\parencite{Pieri2018Modelling,Kancs2016R&D,Onyekwere2019The,Parisi2006Productivity}. This finding points to a modest but robust effect. 

Research in emerging economies including India, Vietnam, China, and Iran yields parallel conclusions \parencite{Manogna2021Does,Zhu2021How,Ding2024R&D,Tung2021The}. Studies indicate that R\&D intensity positively influences sales growth, profitability, returns on assets and equity, and TFP, often outperforming the impact of fixed capital investment. A key mechanism is that R\&D raises the likelihood of both product and process innovation, which subsequently drives productivity gains.

The relationship, however, is not uniformly linear \parencite{Kancs2016R&D}. Returns to R\&D are frequently characterized by threshold effects: low levels of R\&D intensity may yield negligible or zero gains, while investments surpassing a certain intensity exhibit a strong, positive impact on productivity. The input-accumulation channel of R\&D is somewhat stronger in high- and medium-high-technology manufacturing than in low-tech industries, although \textcite{Pieri2018Modelling} document that the technical-efficiency channel runs in the opposite direction (R\&D raises efficiency in low-tech sectors and lowers it in high-tech sectors that operate close to the frontier). Some studies, particularly in specific industries or developing-country contexts, report mixed or small coefficients, highlighting persistent challenges related to endogeneity, measurement error, and methodological constraints \parencite{Onyekwere2019The,Zhu2021How}. Moreover, the productivity payoff from R\&D is often contingent on complementary assets, such as information and communication technologies (ICT), organizational capital, export orientation, and favorable regulatory frameworks.

The following paragraphs explores how these relationships manifest in the Russian market, synthesizing results from relevant empirical studies. Applications of the \textcite{crepon1998research} model to Russian data highlight its utility in addressing transitional economy challenges. For instance, \textcite{mariev2022effects} applied the \textcite{crepon1998research} framework to cross-sectional firm-level data from the Business Environment and Enterprise Performance Survey (BEEPS) for 2012–2014, supplemented with a regional innovation index, focusing on manufacturing firms with varying technological intensities. Their findings confirm that R\&D expenditures, alongside personnel training and cooperation with universities, positively influence innovative sales, which in turn boost productivity, with stronger effects observed in high-tech firms. While the results seem reasonable methods of estimation and dataset underscores the \textcite{crepon1998research} model's strength. 

In contrast, implementations of the \textcite{griliches1979issues} knowledge capital model on Russian data implement panel structures and extensions to handle endogeneity, emphasizing direct augmentation of production functions with indicators of technological progress. \textcite{malein2019total} estimated augmented Cobb-Douglas production functions on panel data for 39 Russian black metal firms from 2006 to 2017, incorporating a technology switch dummy for the transition to electric arc furnaces. Authors found that adopting the new technology increases TFP by 24–28\% and alters production function parameters, reducing labor and material intensities. This single-equation approach, while susceptible to simultaneity biases without such controls, effectively captures level effects of technological change on productivity, aligning with \textcite{griliches1979issues} focus on knowledge-driven shifts but adapted to sector-specific innovations in a transitional context.

In summary, across diverse geographic and economic contexts, R\&D investment is consistently linked to higher productivity and sustained long-term performance. The magnitude of this effect is maximized when R\&D is conducted at sufficient intensity, integrated with complementary intangible assets, and embedded within supportive institutional and innovation ecosystems.



\section{Methods}

\subsection{Models}
\label{subsec:methods_models}

In terms of the empirical framework employed in this study rests on the Cobb-Douglas production function, which relates a firm's output to its primary physical inputs-labour, capital, and intermediate materials. In logarithmic form, the function yields a linear specification in which the estimated coefficients represent the output elasticities of these conventional inputs. Critically, total factor productivity (TFP) is not directly observed but is recovered as a residual: the component of output variation that remains after accounting for the measured contributions of labour, capital, and materials. In all specifications considered below, TFP, denoted $\omega_{it}$, is thus obtained as the difference between actual output and the estimated input contributions, and it is this residual productivity term that constitutes the central object of analysis.

A tiered modeling strategy is employed to estimate the relationship between innovation inputs and firm-level productivity. Notably, the empirical approach advances from reduced-form specifications, which serve to establish baseline correlations, to structural econometric models designed to address endogeneity, selection bias, and the dynamic nature of technological change. This progression permits a rigorous assessment of how the estimated returns to R\&D are influenced by methodological sophistication.

\paragraph{Baseline Reduced-Form Specification}

The foundational approach is grounded in the knowledge capital framework established by \textcite{griliches1979issues}. This framework expands upon the standard Cobb-Douglas production function by treating R\&D as a distinct, independent factor of production. Hence, the core intuition posits that a firm's output is determined not only by conventional inputs (labor, capital, and materials) but also by an accumulated stock of knowledge derived from historical R\&D investments. Overall, this supplemented relationship is expressed as:

\begin{equation}
y_{it} = \beta_0 + \beta_l l_{it} + \beta_k k_{it} + \beta_m m_{it} + \beta_r r_{it} + \mu_{it},
\end{equation}

\noindent where $y_{it}$, $l_{it}$, $k_{it}$, and $m_{it}$ denote the natural logarithms of real revenue, labor, capital, and intermediate inputs, respectively. The variable $r_{it} = \ln(S_{it}+1)$ is the log R\&D capital stock, where $S_{it}$ is the perpetual-inventory R\&D stock defined in Subsection Variables Construction, and $\mu_{it}$ is the composite error term. The coefficient $\beta_r$ captures the output elasticity of R\&D capital and replaces the symbol $\gamma$ used in the original \textcite{griliches1979issues} formulation; the same notation is retained throughout the rest of the thesis for consistency. To account for unobserved heterogeneity, the error term is decomposed as $\mu_{it} = \alpha_i + \lambda_t + \epsilon_{it}$, where $\alpha_i$ represents time-invariant firm-specific effects, $\lambda_t$ captures common time shocks, and $\epsilon_{it}$ is the idiosyncratic error.

Identification in this baseline model relies on the strict exogeneity of input choices relative to the error term, formally stated as $E[\epsilon_{it} | l_{it}, k_{it}, m_{it}, r_{it}] = 0$. Typically, estimation is performed via Ordinary Least Squares (OLS) or Fixed Effects (FE) using the within-transformation to eliminate $\alpha_i$. The primary advantage of this specification is its transparency and comparability with existing literature \parencite{shao2020environmental}. However, this framework is subject to several well-documented sources of bias. First, simultaneity bias arises because productive firms are likely to invest more in R\&D and employ more inputs concurrently, inducing an upward bias in $\beta_r$. Second, omitted variable bias may result from the failure to control for unobserved firm-specific characteristics, such as managerial quality or strategic orientation, that are correlated with both R\&D intensity and output. Third, selection bias is introduced if firm survival is correlated with productivity, such that the estimation sample is non-randomly truncated. Finally, this static framework ignores the dynamic accumulation of knowledge and the persistence of productivity shocks, limiting its causal interpretability.


\paragraph{Structural Control Function Approaches}

To address the simultaneity bias inherent in the reduced-form models, structural control function approaches are employed. These methods utilise the dynamic optimisation behaviour of firms to proxy for unobserved productivity. The general production function is written as: 

\begin{equation}
y_{it} = \beta_l l_{it} + \beta_k k_{it} + \beta_m m_{it} + \beta_r r_{it} + \omega_{it} + \epsilon_{it}
\end{equation}

, where $\omega_{it}$ is the unobserved productivity term known to the firm but not the researcher. R\&D capital is treated as a predetermined state variable analogous to physical capital, determined by lagged information and thus invariant to contemporaneous productivity shocks. Accordingly, $\beta_r$ is identified jointly with $\beta_k$ via the second-stage moment conditions and is subsumed within the first-stage composite $\phi_t$, following the structural treatment of physical capital.

The \textcite{olley1992dynamics} approach proposes using investment as a proxy for unobserved productivity, relying on the intuition that investment is a monotonic function of productivity conditional on the state variables. In terms of the investment demand function is specified as $i_{it} = i_t(\omega_{it}, k_{it}, r_{it})$, where $r_{it}$ denotes the log R\&D capital stock. Under the assumption of monotonicity, this function can be inverted to express productivity as $\omega_{it} = h_t(i_{it}, k_{it}, r_{it})$. Substituting this into the production function yields the first-stage estimation equation:
\begin{equation}
y_{it} = \beta_l l_{it} + \beta_m m_{it} + \phi_t(i_{it}, k_{it}, r_{it}) + \epsilon_{it},
\end{equation}
\noindent where $\phi_t(i_{it}, k_{it}, r_{it}) = \beta_k k_{it} + \beta_r r_{it} + h_t(i_{it}, k_{it}, r_{it})$ is approximated by a higher-order polynomial. Labour and materials enter as free variables, chosen after the productivity shock is realized, and their coefficients are identified at this stage, while the state-variable coefficients (capital and knowledge stock) are absorbed into the composite function $\phi_t(\cdot)$ and identified in a subsequent stage. This method assumes investment is strictly positive and monotonically increasing in productivity. However, in the Russian context, investment data is often lumpy or zero-inflated due to accounting practices, potentially violating the monotonicity condition and leading to sample selection bias.

Given the limitations of investment data, the primary reliance is placed on the \textcite{levinsohn2003estimating} approach, which uses intermediate inputs as the proxy. The intuition here is that intermediate inputs, such as materials and energy, respond more flexibly to productivity shocks than capital or investment. Productivity is recovered via the demand for intermediates, specified as $m_{it} = m_t(\omega_{it}, k_{it})$, which is inverted to $\omega_{it} = f^{-1}(m_{it}, k_{it})$.

In the original \textcite{levinsohn2003estimating} formulation, the production function is estimated in two stages: in the first stage, labour and material coefficients are recovered from the regression $y_{it} = \beta_l l_{it} + \phi_t(m_{it}, k_{it}) + \epsilon_{it}$, and in the second stage, the capital coefficient is identified using the Markov evolution of productivity. However, as demonstrated by \textcite{ackerberg2015identification}, this procedure suffers from a functional dependence problem. Because labour is typically chosen simultaneously with materials and both are functions of the same state variables (capital and productivity), the first-stage regression can not separately identify $\beta_l$ from the nonparametric proxy function $\phi_t(\cdot)$. To overcome this identification hurdle, the correction proposed by \textcite{ackerberg2015identification} is adopted. In the first stage, the composite function $\phi_t(l_{it}, m_{it}, k_{it}, r_{it})$ is estimated nonparametrically via a higher-order polynomial in all observed inputs (labour, materials, capital, and the R\&D capital stock $r_{it}$) without attempting to recover any individual input coefficients. This stage serves exclusively to eliminate measurement error $\epsilon_{it}$ from the production function and to construct the predicted composite $\hat{\phi}_{it}$. In the second stage, all four production function coefficients $(\beta_l, \beta_m, \beta_k, \beta_r)$ are identified simultaneously via two-step GMM, with R\&D capital treated symmetrically with physical capital as a state variable rather than absorbed into the proxy polynomial. The productivity term is recovered as $\hat{\omega}_{it} = \hat{\phi}_{it} - \beta_l l_{it} - \beta_m m_{it} - \beta_k k_{it} - \beta_r r_{it}$, and its Markov evolution is specified as $\omega_{it} = \rho \omega_{it-1} + \xi_{it}$, where $\xi_{it}$ is the innovation shock orthogonal to information available at $t-1$. The coefficient vector $(\beta_l, \beta_m, \beta_k, \beta_r)$ is then estimated by minimising the sample analogue of the moment conditions $E[\xi_{it}(\beta) \cdot Z_{it}] = 0$, where $Z_{it}$ stacks lagged labour, lagged materials, lagged capital, lagged R\&D stock, current capital, and current R\&D stock are six instruments for four parameters, yielding two over-identifying restrictions tested by the Hansen $J$ statistic.

\paragraph{Why the LP-ACF specification with R\&D as a state variable works.} Identifying $\beta_r$ jointly with $\beta_k$ in Stage~2 is methodologically sound under three conditions, all of which are met in the present sample. First, R\&D capital evolves as a perpetual-inventory stock with depreciation $\delta = 0.15$ (Section~Variables Construction), so $r_{it}$ is determined by past R\&D expenditure choices and the depreciation law alone, not by today's productivity shock $\xi_{it}$; this is the textbook state-variable property that justifies using $r_{it}$ itself (and its lag) as instruments. Second, the materials-as-proxy inversion that recovers $\omega_{it}$ already conditions on $r_{it}$ through the polynomial $\phi_t(l, m, k, r)$, so the productivity series fed into the Markov transition is conditional on the firm's R\&D stock, exactly as it must be for the lagged moment conditions to be valid \parencite{ackerberg2015identification}. Third, the order condition is satisfied (six moments versus four parameters), and the rank condition is supported empirically by the substantial within-firm autocorrelation of $r_{it}$ (a consequence of the perpetual-inventory construction and the persistence of innovator status documented in the descriptive analysis), which gives the lagged R\&D stock genuine identifying power for $\beta_r$. Treating R\&D only as a polynomial argument absorbed into $\phi_t$, which was the standard LP-ACF practice prior to the modification adopted here, left $\beta_r$ inseparable from $\omega_{it}$ in the second stage and forced reliance on OP and DJ alone for the R\&D elasticity. The current specification removes that asymmetry and allows LP to deliver an estimate of $\beta_r$ that is directly comparable to the OP point estimate while sharing the LP attractive feature of preserving firms with zero investment.

Ultimately, the primary advantage of the \textcite{levinsohn2003estimating} approach combined with the \textcite{ackerberg2015identification} correction is threefold: it retains firms with zero investment, preserving sample size crucial for the RFSD panel \parencite{bondarkov2025russian}; it ensures that all input elasticities, now including the R\&D capital elasticity $\beta_r$, are consistently identified under the same set of moment conditions; and the over-identifying restrictions it places on the model can be formally tested via the Hansen $J$ statistic. A potential limitation is the assumption that intermediate inputs are chosen after productivity is realised but before output is produced; bias may arise if input prices vary across firms and are correlated with productivity.

\paragraph{Structural Dynamic Innovation Model}

The most advanced specification in this research adopts the framework of \textcite{doraszelski2013r}. Unlike the \textcite{olley1992dynamics}  and \textcite{levinsohn2003estimating} frameworks which treat productivity as evolving exogenously via a Markov process independent of firm actions, this model endogenizes the productivity transition process through R\&D. Subsequently, the central assumption here is that R\&D does not simply enter the production equation as a fixed variable; instead, it serves as the mechanism by which the very landscape of future productivity, especially its probabilistic outcomes, is recalibrated. The evolution of productivity $\omega_{it}$ is modeled as a controlled Markov process:

\begin{equation}
\omega_{it} = g(\omega_{it-1}, r_{it-1}) + \xi_{it},
\end{equation}

\noindent where $r_{it-1}$ denotes R\&D expenditures in the previous period, $g(\cdot)$ is the transition function determining expected efficiency gains, and $\xi_{it}$ is an innovation shock orthogonal to information at $t-1$. Following \textcite{doraszelski2013r}, the transition function is approximated using a complete polynomial of degree three in $(\omega_{it-1}, r_{it-1})$:
\begin{equation}
\begin{split}
\omega_{it} = {} & \gamma_0 + \gamma_1 \omega_{it-1} + \gamma_2 r_{it-1}
 + \gamma_3 \omega_{it-1}^2 + \gamma_4 r_{it-1}^2
 + \gamma_5 \omega_{it-1} r_{it-1} \\
 & + \gamma_6 \omega_{it-1}^3 + \gamma_7 r_{it-1}^3
 + \gamma_8 \omega_{it-1}^2 r_{it-1}
 + \gamma_9 \omega_{it-1} r_{it-1}^2
 + \xi_{it}.
\end{split}
\label{eq:dj_transition}
\end{equation}

To summarize, the primary advantage of this approach is that it provides a causal estimate of how innovation investments reshape the firm's long-term efficiency frontier, capturing the dynamic option value of R\&D. However, the model is computationally intensive, requires strong assumptions about the functional form of the productivity transition, and is sensitive to measurement error in R\&D, which is a known issue in Russian accounting data. Thus, the hierarchical modeling strategy ensures that baseline correlations are established, externalities are accounted for, standard simultaneity bias is corrected via control functions, and finally, the dynamic causal mechanism of innovation on productivity growth is isolated.


\subsection{Estimation Strategy}

In terms of the estimation strategy proceeds hierarchically, aligning with the model specification. Initial estimates are obtained using standard panel data techniques to establish baseline correlations. Three structural estimators are then applied sequentially: the \textcite{olley1992dynamics} control function approach with survival correction, the \textcite{levinsohn2003estimating} control function approach with the \textcite{ackerberg2015identification} identification correction, and the \textcite{doraszelski2013r} dynamic innovation model. Indeed, this multi-stage approach ensures that the robustness of the findings is tested against increasingly stringent identification assumptions.

\paragraph{Identification logic of the hierarchy.} Each layer of the hierarchy isolates a different source of identification failure rather than producing an alternative answer to the same question. Pooled OLS recovers the descriptive cross-sectional association between R\&D and output but conflates simultaneity, unobserved heterogeneity, and the endogeneity of the productivity transition. To differentiate between selection effects and internal causal links, we employ within-firm fixed effects to partial out time-invariant characteristics like managerial quality and vintage, effectively quantifying how much of the OLS R\&D coefficient is attributable to firm-specific selection. The proxy-variable estimators of \textcite{olley1992dynamics} and \textcite{levinsohn2003estimating} address the residual simultaneity between variable inputs and the productivity shock $\xi_{it}$ by inverting an investment or materials-demand function for $\omega_{it}$, but they share with OLS and FE the assumption that $\omega_{it}$ evolves according to an \emph{exogenous} first-order Markov law $\omega_{it} = \rho \omega_{it-1} + \xi_{it}$ that is independent of firm actions. It is exactly this last assumption that the \textcite{doraszelski2013r} specification relaxes: the transition function $g(\omega_{it-1}, r_{it-1})$ is allowed to depend parametrically on lagged R\&D, so innovation becomes an endogenous control over the productivity law of motion rather than a passive input. All four structural estimators (FE, OP, LP, DJ) thus impose a Markov structure on $\omega_{it}$, the genuine distinction is whether $r_{it}$ is permitted to shift the conditional distribution of next-period productivity. The hierarchy therefore decomposes the OLS--DJ gap into three economically meaningful components: a selection component absorbed by FE, a contemporaneous-simultaneity component absorbed by OP/LP, and an endogenous-transition component identified only by DJ. This decomposition motivates the empirical contribution of the thesis: rather than pitting a `correct' estimator against a `naive' one, the goal is to quantify how much of the static R\&D-output association reflects each of these three identification problems on Russian production-economy data.

%%  REDUCED-FORM SPECIFICATIONS

\paragraph{Reduced-Form Specifications}

The baseline reduced-form specifications are estimated via pooled OLS with \textcite{mackinnon1985some} heteroscedasticity-consistent standard errors - the bias-correction variant that reweights each squared residual by $1/(1-h_{ii})^2$, where $h_{ii}$ is the leverage of observation $i$. Both Fixed Effects and Two-Way Fixed Effects (entity and time) specifications employ firm-level cluster-robust standard errors to account for within-firm heteroscedasticity and serial correlation. Because these estimators control only for time-invariant heterogeneity and not for time-varying simultaneity bias, the results are interpreted as descriptive benchmarks rather than causal estimates.

%%  OLLEY–PAKES

\paragraph{Olley-Pakes Control Function Estimator}

The \textcite{olley1992dynamics} control function estimator is implemented as a three-stage procedure that addresses both simultaneity bias and survival bias in a Cobb-Douglas production function framework.

\paragraph{Stage~1: Identification of free-variable coefficients.}
The key identification assumption is that investment is a strictly monotone function of unobserved productivity conditional on the state variables: $i_{it} = f_t(\omega_{it}, k_{it}, r_{it})$, where $i_{it}$ denotes investment, $k_{it}$ denotes capital, and $r_{it}$ the knowledge stock. Inverting this demand yields $\omega_{it} = h_t(i_{it}, k_{it}, r_{it})$, which is substituted into the production function:
\begin{equation}
  y_{it}
    = \beta_l \, l_{it} + \beta_m \, m_{it}
      + \underbrace{\beta_0 + \beta_k \, k_{it} + \beta_r \, r_{it}
        + h_t(i_{it}, k_{it}, r_{it})}_{\displaystyle \phi_t(k_{it}, r_{it}, i_{it})}
      + \epsilon_{it},
  \label{eq:op_s1}
\end{equation}
where the composite $\phi_t(\cdot)$ is approximated by a third-degree polynomial in $(k_{it}, r_{it}, i_{it})$. Because labour and materials are \emph{free} variables, chosen after the productivity shock is realized, they enter linearly and their coefficients $\hat\beta_l$ and $\hat\beta_m$ are identified at this stage by OLS. The predicted composite $\hat\phi_{it}$ is retained for Stage~2. The monotonicity assumption restricts the estimation sample to firm-year observations with strictly positive investment ($i_{it} > 0$).

\paragraph{Stage~2: Survival correction.}
To control for non-random exit, a probit model of firm survival is estimated using a second-degree polynomial in the lagged state variables:
\begin{equation}
  \Pr(\text{survive}_{it} = 1)
    = \Phi\!\bigl(\delta_0 + \delta_1 k_{it-1} + \delta_2 r_{it-1}
      + \delta_3 k_{it-1}^2 + \delta_4 r_{it-1}^2
      + \delta_5 k_{it-1} r_{it-1}\bigr),
  \label{eq:op_probit}
\end{equation}
where $\Phi(\cdot)$ denotes the standard normal CDF. Following the additive-index variant of the \textcite{olley1992dynamics} correction, the predicted survival probability $\hat{P}_{it}$ is included as an \emph{additional state-variable control} inside the Stage-3 productivity transition polynomial $g(\omega_{it-1}, \hat{P}_{it})$. Intuitively, $\hat{P}_{it}$ summarises how the distribution of future productivity conditional on the state differs for surviving versus exiting firms, so augmenting $g(\cdot)$ with $\hat{P}_{it}$ absorbs the selection wedge in the transition function without resorting to an explicit inverse-Mills or inverse-probability-weighting correction. This is the implementation used in the estimation code and mirrors the treatment in \textcite{olley1992dynamics}.

\paragraph{Stage~3: Identification of state-variable coefficients.}
Define $\hat\omega_{it}(\beta_k, \beta_r) \equiv \hat\phi_{it} - \beta_k k_{it} - \beta_r r_{it}$. The Markov transition of productivity conditional on survival is approximated as
\begin{equation}
  \omega_{it} = g(\omega_{it-1},\, \hat{P}_{it}) + \xi_{it},
  \label{eq:op_transition}
\end{equation}
where $g(\cdot)$ is a flexible polynomial:
\begin{equation}
  g(\omega_{it-1}, \hat{P}_{it})
    = \gamma_0 + \gamma_1 \omega_{it-1} + \gamma_2 \omega_{it-1}^2
      + \gamma_3 \omega_{it-1}^3 + \gamma_4 \hat{P}_{it}
      + \gamma_5 \hat{P}_{it}^2 + \gamma_6 \omega_{it-1} \hat{P}_{it}.
\end{equation}
For each candidate $(\beta_k, \beta_r)$, the $\gamma$ vector is profiled analytically via OLS, and the innovation $\xi_{it}$ is computed as the residual. The state-variable coefficients are then obtained by nonlinear least squares (NLS):
\begin{equation}
  (\hat\beta_k,\, \hat\beta_r)
    = \arg\min_{\beta_k,\, \beta_r}
      \sum_{i,t} \xi_{it}(\beta_k, \beta_r)^2,
  \label{eq:op_nls}
\end{equation}
with numerical optimization performed via the Nelder-Mead simplex algorithm. This procedure yields TFP estimates $\hat\omega_{it}^{OP} = \hat\phi_{it} - \hat\beta_k k_{it} - \hat\beta_r r_{it}$ that are purged of simultaneity bias (via the control function), survival bias (via the probit correction), and measurement error (via the first-stage polynomial).

Taken together, the classical \textcite{olley1992dynamics} procedure identifies free-variable coefficients in Stage 1 by assuming they are chosen after observing $\omega_{it}$ but are conditionally independent of the polynomial control once the state variables and the proxy are included. \textcite{ackerberg2015identification} show that this assumption may fail when labour and materials are collinear with the polynomial terms, creating a functional dependence problem. The OP estimates are therefore reported alongside the ACF-corrected LP estimates as a robustness benchmark. The LP approach also has the practical advantage of preserving firms with zero investment \parencite{bondarkov2025russian}, which are excluded from the OP sample by the monotonicity requirement on $i_{it} > 0$.

%%  LEVINSOHN–PETRIN WITH ACF CORRECTION

\paragraph{Levinsohn-Petrin Estimator with ACF Correction}

To address the functional dependence problem identified by \textcite{ackerberg2015identification}, the \textcite{levinsohn2003estimating} control function approach is implemented with the ACF identification correction. Hence, materials expenditure serves as the proxy for unobserved productivity, replacing investment and thereby preserving firms with zero investment in the estimation sample.

\paragraph{Stage~1: Composite function estimation.}
The first-stage specification is:
\begin{equation}
y_{it} = \phi_t(l_{it}, m_{it}, k_{it}, r_{it}) + \epsilon_{it},
\end{equation}
where $\phi_t(\cdot)$ is approximated by a third-degree polynomial in all observed inputs, including labour, materials, capital, and the knowledge stock. This composite function absorbs all production function coefficients and the unobserved productivity term; no individual input elasticities are identified at this stage, in contrast to the classical LP procedure. Hence, the first stage serves exclusively to purge measurement error and to recover the predicted composite $\hat{\phi}_{it}$.

\paragraph{Stage~2: Two-step GMM.}
In the second stage, the full coefficient vector $(\beta_l, \beta_m, \beta_k)$ is identified jointly via two-step GMM. For a candidate parameter vector $\beta$, the productivity term is constructed as
\begin{equation}
\hat{\omega}_{it}(\beta) = \hat{\phi}_{it} - \beta_l l_{it} - \beta_m m_{it} - \beta_k k_{it},
\end{equation}
and its Markovian evolution $\omega_{it} = \rho_0 + \rho_1 \omega_{it-1} + \rho_2 \omega_{it-1}^2 + \rho_3 \omega_{it-1}^3 + \xi_{it}$ is estimated via OLS to extract the innovation shocks $\xi_{it}$. The coefficient vector is then estimated by minimizing the sample analogue of the moment conditions $E[\xi_{it}(\beta) \cdot Z_{it}] = 0$.

The instrument vector $Z_{it}$ exploits the timing assumptions of the model:
\begin{equation}
Z_{it} = (l_{it-1},\; m_{it-1},\; k_{it-1},\; k_{it})'.
\end{equation}
Lagged labour and materials are valid instruments because they were chosen in period $t-1$, before the realization of the innovation shock $\xi_{it}$. Current capital $k_{it}$ is included because it is a state variable determined by investment decisions in $t-1$ and thus predetermined with respect to $\xi_{it}$. This yields $q = 4$ moment conditions for $p = 3$ parameters, providing $q - p = 1$ overidentifying restriction.

Regarding the two-step procedure proceeds as follows. In the first step, the identity matrix serves as the weighting matrix: $W_1 = I_q$. Using the first-step residuals $\hat\xi_{it}^{(1)}$, the optimal weighting matrix is constructed as
\begin{equation}
\hat{W} = \left[\frac{1}{N} \sum_{i,t} \hat\xi_{it}^{(1)2}\, Z_{it} Z_{it}'\right]^{-1},
\end{equation}
which accounts for heteroscedasticity in the innovation process. The second step re-estimates $\beta$ using this optimal weighting matrix. Numerical optimization employs the L-BFGS-B algorithm with bound constraints to ensure economically meaningful parameter values.

%%  DORASZELSKI–JAUMANDREU

The structural dynamic innovation model of \textcite{doraszelski2013r} generalizes the Markov productivity process by allowing R\&D expenditure to endogenously shift the distribution of future productivity. The first stage is identical to the ACF procedure: a third-degree polynomial in $(l_{it}, m_{it}, k_{it}, r_{it})$ yields the composite $\hat\phi_{it}$, purging measurement error without identifying individual input elasticities.

\paragraph{Transition function.}
In the second stage, the productivity transition equation is specified as a complete third-degree polynomial in lagged productivity and R\&D expenditure, following the original specification of \textcite{doraszelski2013r} in \eqref{eq:dj_transition}. The ten parameters of the cubic polynomial capture non-linear persistence ($\gamma_3$, $\gamma_6$), diminishing returns to R\&D ($\gamma_4$, $\gamma_7$), and heterogeneous complementarity between productivity and R\&D through the cross-derivative $\partial^2 g / \partial\omega\partial r = \gamma_5 + 2\gamma_8\omega_{it-1} + 2\gamma_9 r_{it-1}$, which is allowed to vary across firms rather than being constrained to a constant. For each candidate vector of production function parameters $\theta = (\beta_l, \beta_m, \beta_k)$, the transition parameters $\gamma$ are profiled analytically via least squares, ensuring computational efficiency.

\paragraph{GMM identification.}
The production function parameters are identified through GMM, exploiting the orthogonality between the innovation shock $\xi_{it}$ and the information set available to the firm at the time of investment. The instrument vector is
\begin{equation}
Z_{it} = (l_{it-1},\; m_{it-1},\; k_{it},\; r_{it-1},\; k_{it-1})'.
\end{equation}
This provides $q = 5$ moment conditions for $p = 3$ production function parameters, yielding $q - p = 2$ overidentifying restrictions. The inclusion of both $k_{it}$ and $k_{it-1}$ is justified by the state-variable timing: capital is predetermined with respect to $\xi_{it}$, while lagged capital provides an additional source of exogenous variation. The parameter vector is estimated by minimizing the quadratic form of the sample moments:
\begin{equation}
\hat{\theta}_{GMM} = \arg \min_{\theta} \left( \frac{1}{NT} \sum_{i,t} \xi_{it}(\theta) Z_{it} \right)' W \left( \frac{1}{NT} \sum_{i,t} \xi_{it}(\theta) Z_{it} \right),
\end{equation}
where $W$ is a positive definite weighting matrix. As with the LP estimator, a two-step procedure is employed: the first step uses $W_1 = I_5$ and the second step uses the heteroscedasticity-consistent optimal weighting matrix computed from first-step residuals.

\paragraph{Hansen $J$-test.}
The validity of the overidentifying restrictions is tested using Hansen's $J$-statistic, distributed as $\chi^2$ with degrees of freedom equal to the number of overidentifying restrictions. For the LP estimator, $J \sim \chi^2(1)$; for the DJ estimator, $J \sim \chi^2(2)$. Given the large sample size ($N > 900{,}000$), the $J$-statistic is liable to reject even negligible misspecification; accordingly, we supplement the $p$-value with the normalized statistic $J/N$ as a measure of the economic magnitude of instrument invalidity. This approach isolates the causal impact of R\&D on the drift of the productivity process, distinct from its level effect.

%.    \parencite{hall2005generalized}.    %
%%  ENDOGENEITY AND INSTRUMENTS

Endogeneity of R\&D investment is addressed across all structural specifications. In the control function approaches (OP and LP-ACF), the simultaneity between input choices and productivity is addressed by the proxy variable itself, which absorbs the contemporaneous productivity shock into the composite $\phi_t(\cdot)$. In the DJ model, endogeneity is moderated by the use of lagged instruments, under the assumption that while firms anticipate future productivity when investing, they cannot predict the realization of the innovation shock $\xi_{it}$.

%%  BOOTSTRAP INFERENCE

Standard errors for the OP, LP, and DJ estimators are computed via a clustered panel bootstrap to account for the generated-regressor problem inherent in multi-stage estimation. The procedure resamples at the \emph{firm} level with replacement, preserving the entire time-series history of each drawn firm, and re-estimates the full multi-stage pipeline on each bootstrap sample. Let $B$ denote the number of bootstrap replications. The bootstrap variance estimator is constructed as:
\begin{equation}
\widehat{Var}_{boot}(\hat{\beta}) = \frac{1}{B-1} \sum_{b=1}^{B} (\hat{\beta}^{(b)} - \bar{\beta})(\hat{\beta}^{(b)} - \bar{\beta})',
\end{equation}
where $\hat{\beta}^{(b)}$ is the estimate from the $b$-th bootstrap sample and $\bar{\beta} = B^{-1}\sum_{b}\hat{\beta}^{(b)}$ is the bootstrap mean. Confidence intervals are constructed using the percentile method: the $100(1-\alpha)\%$ interval for $\beta_j$ is given by the $(\alpha/2)$ and $(1-\alpha/2)$ quantiles of $\{\hat{\beta}_j^{(b)}\}_{b=1}^{B}$. $p$-values are obtained as
\begin{equation}
p_j = 2 \cdot \min\!\left(
    \frac{1}{B}\sum_{b=1}^{B}\mathbf{1}[\hat{\beta}_j^{(b)} < 0],\;
    \frac{1}{B}\sum_{b=1}^{B}\mathbf{1}[\hat{\beta}_j^{(b)} > 0]
\right),
\end{equation}
which provides a distribution-free analogue of a two-sided test. We use $B = 200$ replications for the baseline results.\footnote{$B = 200$ is standard for panel-bootstrap inference in applied production-function estimation; Monte-Carlo noise in the resulting standard errors is small relative to their analytical magnitude on a panel of this size ($N \approx 691{,}644$ firm-year observations).}

This bootstrap strategy is necessitated by three considerations: (i)~the generated-regressor problem, whereby the use of first-stage estimated productivity $\hat{\phi}_{it}$ in the second-stage GMM renders conventional analytic standard errors invalid; (ii)~the nonlinearity of the transition function, particularly in the DJ model where polynomial interactions between lagged productivity and R\&D preclude closed-form variance expressions; and (iii)~the need to preserve within-firm temporal dependence, which is maintained by resampling at the firm (panel) level rather than at the observation level.

%%  ROBUSTNESS CHECKS

Several robustness checks are implemented to ensure the reliability of the estimates. First, the sensitivity of the results to outlier influence is assessed by trimming all continuous variables at the 1st and 99.5th percentiles. Specifically, for each variable $x$, observations satisfying $x < P_{1}$ or $x > P_{99.5}$ are excluded from the estimation sample, where $P_{k}$ denotes the $k$-th percentile of the empirical distribution. Second, the sensitivity of the structural estimators to the polynomial degree of the control function is assessed by re-estimating the OP and LP models with second-, third-, and fourth-degree polynomials. The stability of the input elasticities across these specifications provides evidence on the adequacy of the nonparametric approximation.

Potential identification threats, such as measurement error in the accounting-based R\&D variable and survival bias from firm exit, are acknowledged. In the OP estimator, survival bias is explicitly controlled through the probit selection equation. In the LP and DJ frameworks, the estimation sample includes exiting firms up to their final year of observation, and the dynamic framework naturally accounts for selection through the Markov transition process. Collectively, these methodological safeguards ensure that the estimated relationship between innovation and productivity reflects structural economic mechanisms rather than statistical artefacts.


\section{Data and Variables}

\subsection{Data Source}

The empirical analysis in this study is based on the Russian Financial Statements Database (RFSD) \parencite{bondarkov2025russian}, a comprehensive and harmonized panel of annual unconsolidated financial statements covering the universe of Russian firms from 2011 to 2023. The RFSD \parencite{bondarkov2025russian} represents a significant advancement in data availability for the Russian market, as it includes the first open-source records of every active firm in the country, thereby mitigating the selection bias often found in proprietary databases that exclude non-filing or smaller entities.

To estimate the production function parameters and the endogenous productivity process discussed in the previous section, we construct a longitudinal panel dataset comprising output, primary inputs, and innovation-specific variables. Output ($Y_{it}$) is measured as total firm revenue, while capital ($K_{it}$) is represented by the book value of fixed assets. For the structural estimation using the \textcite{levinsohn2003estimating} framework, we utilize intermediate consumption ($M_{it}$), defined as the total cost of raw materials, energy, and fuel. Where the \textcite{olley1992dynamics} inversion is preferred, annual investment ($I_{it}$) is used, proxied by the change in fixed assets adjusted for depreciation. A critical choice in the research is the measurement of labor ($L_{it}$). While traditional studies often rely on headcount, the reported number of employees in Russian corporate registries (such as SPARK) is frequently fragmented, characterized by a high incidence of missing values and zeros. Robustness of the estimations and the need to capture heterogeneity in labor quality motivate the use of total wage expenditures as a proxy for labor services. As argued previously, this value-based approach effectively weights labor units by their market-determined marginal productivity, thereby adjusting for differences in human capital across firms.

The primary independent variable of interest is the investment in Research and Development ($R_{it}$). Within the Russian accounting framework, R\&D expenditures are conventionally reported in line 1120 of the balance sheet. However, utilizing this variable presents a non-trivial identification challenge. Historically, only a relatively small subset of enterprises, approximately 4,000 firms annually, explicitly populate this line. This sparsity is partly due to accounting discretion; firms may capitalize R\&D as intangible assets or write them off as current period expenses. Indeed, following regulatory shifts toward the elimination of line 1120 in 2024, most R\&D activities are expected to be subsumed under intangible assets.

To address this, the panel structure of the RFSD \parencite{bondarkov2025russian} is leveraged to track R\&D intensity over time, though the potential under-reporting of innovation activities is acknowledged as a limitation. All monetary variables are subsequently deflated to constant 2011 prices to ensure comparability in real terms across the 2011-2023 observation window, as elaborated below.

The raw data were extracted year by year from the HuggingFace repository, selecting a targeted set of variables: the firm identifier (INN), the regional geographic code (OKTMO), the reporting year, the industry classification at the OKVED section level, revenue (balance sheet line 2110), fixed assets (line 1150), cost of sales (line 2120), research and development expenditures (line 1120), and cash flows attributable to labour (line 4122). At the point of extraction, observations flagged as outliers by the source dataset's internal quality controls were removed, and rows containing null values in any key variable other than R\&D (line 1120) were dropped. After consolidating all annual files into a unified longitudinal panel, duplicate firm-year records were eliminated by retaining the first occurrence for each unique (INN, year) pair, thereby ensuring a strictly balanced identification structure.

Several initial transformations were then applied to harmonize the raw accounting data for econometric use. Missing R\&D values were imputed with zeros under the assumption that the absence of a reported figure on line 1120 indicates no formal R\&D activity in the given year. Because Russian accounting conventions record cash outflows as negative values, both the labour cash flow variable ($CFo\_labor$) and the cost of sales ($PL\_cost\_of\_sales$) were multiplied by $-1$ to yield positive figures amenable to logarithmic transformation. Industry labels were mapped from OKVED section codes to their full-text descriptions to facilitate sector-level analysis.

To ensure comparability across years, all monetary variables were deflated to constant 2011 prices using the Consumer Price Index (CPI). The deflation procedure distinguishes between flow and stock variables, reflecting their different temporal measurement conventions. For flow variables, which represent cumulative activity over the fiscal year (revenue, cost of sales, labour expenditure), a deflator based on the average annual price level was computed. Specifically, monthly CPI chain indices were cumulated multiplicatively from January 2011, and the resulting monthly price levels were averaged within each calendar year; the ratio of each year's average price level to the 2011 average yields the flow deflator. For balance sheet stock variables (fixed assets, R\&D expenditures), which are measured at year-end, the deflator was constructed from the cumulative end-of-December price level relative to the 2011 annual average. This dual-deflator approach avoids the temporal mismatch that would arise from applying a single price index to variables measured at different points within the fiscal year.

Notably, the selection of the eight specified macroeconomic sectors is predicated on a dual imperative of methodological rigor and structural heterogeneity. From an empirical standpoint, the primary inclusion criterion was robust data sufficiency; specifically, only broad OKVED sections yielding in excess of 1,000 observations were retained to ensure the statistical validity and consistency of production function estimates. 
Beyond this methodological constraint, this specific cross-section such as: spanning primary extraction (e.g., Mining, Agriculture), secondary industry (e.g., Manufacturing, Construction), and critical infrastructure and services (e.g., Information and Communication, Transportation), was strategically chosen to encapsulate a comprehensive spectrum of aggregate economic activity. 
Furthermore, the chosen industries possess a distinct production orientation, allowing for a meaningful parameterization of their output using the Cobb-Douglas framework. 
Service, financial, and public-sector activities are excluded on the grounds that the Cobb--Douglas framework provides a poor structural fit for these sectors, as confirmed by boundary coefficient estimates and implausible returns to scale in preliminary runs. This focused approach serves a dual function: it pinpoints the specific sectors in which R\&D has the most significant impact on production, while simultaneously acting as a robustness check for the stability of the input elasticities recovered in the pooled models.
By encompassing industries with markedly distinct baseline factor elasticities and returns to scale (RTS), this selection provides a robust foundation for a nuanced, comparative analysis of how R\&D investments differentially impact productivity across structurally diverse economic environments.

The full multi-criteria cleaning protocol, the sectoral filter that retains only the eight production-oriented OKVED sections, and the supplementary post-cleaning bivariate diagnostics that informed the calibration of the cleaning thresholds are reported in Appendix~\ref{app:cleaning}.\hypertarget{target:sector_rd_shares_detailed}{}\hypertarget{target:median_log_distributions}{}\hypertarget{target:rd_investment_revenue_comparison}{}\hypertarget{target:revenue_distribution_dynamics}{}\hypertarget{target:spearman_correlation}{}\hypertarget{target:box_rd_assets_comparison}{}\hypertarget{target:dist_gross_profit}{}\hypertarget{target:dist_operating_profit}{}\hypertarget{target:revenue_by_sector_box}{}\hypertarget{target:scatter_rev_rd_assets}{}

\subsection{Variables Construction}
\label{subsec:variables_construction}

Following the cleaning and trimming stages, the data were subjected to a final preparation pipeline prior to estimation. First, a validation step confirmed that all observations retained strictly positive values for revenue, fixed assets, cost of sales, and labour, and non-negative values for R\&D expenditure, with any residual violations removed. The core production variables were then transformed into natural logarithms: $l_y = \ln(\text{Revenue})$, $l_k = \ln(\text{Fixed Assets})$, $l_l = \ln(\text{Labour})$, and $l_m = \ln(\text{Cost of Sales})$. For R\&D, the transformation $\log(R_t + 1)$ was applied to both the annual flow ($l_{r\_inv}$) and the accumulated stock ($l_r$), accommodating the large mass of zero values whilst preserving a monotonic mapping.

The cumulative R\&D capital stock ($S_t$) was constructed using the perpetual inventory method (PIM). Because the actual rate of knowledge obsolescence is unobservable, the annual depreciation rate was set to 15\%:
\begin{equation}
    S_t = S_{t-1} \times (1 - 0.15) + R_t.
\end{equation}

This 15\% rate is the widely accepted benchmark in the innovation literature \parencite{griliches1979issues}, reflecting the assumption that technological knowledge loses its competitive relevance at a constant, relatively rapid pace. The accumulation was initialized at each firm's first observed R\&D expenditure and was reset whenever the panel exhibited a gap in consecutive years for a given firm, thereby ensuring that the stock variable reflects only verifiable, temporally continuous knowledge accumulation. A binary indicator for innovator status ($is\_innovator$) was also generated, taking the value of one whenever a firm reports positive R\&D expenditure in a given year and zero otherwise.

In contrast to the assumed depreciation rate for unobservable knowledge capital, physical capital investment ($I_t = K_t - K_{t-1}$) was calculated directly as the first difference in fixed assets across consecutive years. A universal depreciation rate (such as 15\%) is not applied here because reported fixed assets ($K_t$) are derived from firm balance sheets. These accounting book values inherently embed asset-specific depreciation schedules (e.g., varying rates for buildings versus machinery) as mandated by accounting standards. Calculating the net change in these book values provides a proxy for investment to facilitate \textcite{olley1992dynamics} estimation. This investment variable was set to missing for non-consecutive year transitions.
 
Temporal lags of all key estimation variables ($l_y$, $l_k$, $l_l$, $l_m$, $l_r$, $l_{r\_inv}$, $l_{inv}$) were constructed strictly within consecutive firm-year pairs; lagged values were set to missing whenever a gap of more than one year was detected in a firm's reporting history, preventing spurious dynamic linkages. Observations with missing values in the core log-transformed variables were dropped at the final stage, yielding the estimation-ready panel. 

\subsection{Descriptive Analysis}

\begin{table}[htbp]
\centering
\caption{Descriptive Statistics of Estimation Variables (Logs)}
\label{tab:desc_log}
\begin{tabular}{lrrrrrrr}
\toprule
Variable & $N$ & Mean & Std & Min & Median & Max & Skew. \\
\midrule
  Output ($\ln Y$) & 691,644 & 11.565 & 2.042 & 5.834 & 11.636 & 17.267 & -0.15 \\
  Labor ($\ln L$) & 691,644 & 9.515 & 2.045 & 4.060 & 9.648 & 14.920 & -0.17 \\
  Capital ($\ln K$) & 691,644 & 4.944 & 2.626 & -2.012 & 5.143 & 11.957 & -0.23 \\
  Materials ($\ln M$) & 691,644 & 11.352 & 2.066 & 5.130 & 11.440 & 17.001 & -0.19 \\
  R\&D Stock ($\ln R$) & 691,644 & 0.113 & 0.783 & 0.000 & 0.000 & 11.723 & 7.75 \\
\bottomrule
\end{tabular}
\begin{flushleft}
\footnotesize
\textit{Note:} Estimation sample ($N = 691,644$ firm-year observations).
Variables in natural logarithms.
\end{flushleft}
\end{table}


The estimation-ready panel comprises $691{,}644$ firm-year observations on $162{,}887$ unique firms, drawn from the 2011-2023 RFSD extract and restricted at the outset to the eight production-oriented OKVED sections on which a Cobb-Douglas framework is structurally appropriate (manufacturing, mining, agriculture, construction, electric energy and gas, water supply and sanitation, information and communication, and transportation). The complete data preparation sequence, from initial extraction and sectoral filtering to the cleaning and temporal-lag architecture, is documented in Appendix~\ref{app:cleaning}. 

Table~\ref{tab:desc_log} report level- and log-form descriptive statistics for the four core production variables and the R\&D capital stock. After log transformation, output, labour, capital, and materials exhibit closely aligned means and medians, with skewness coefficients in the narrow band $[-0.23,-0.15]$, comfortably within the parametric tolerance of the Cobb-Douglas specification. Capital displays the widest dispersion ($\sigma_{\ln K} = 2.626$), reflecting the heterogeneity of physical-capital intensities preserved by the cleaning protocol, while output, labour, and materials cluster around comparable dispersions ($\sigma \in [2.04, 2.07]$). Every input retains a strictly positive support after the log transformation, with finite first and second moments and no degenerate mass in the bulk of the distribution; consequently, the four variables are well posed for the OLS, fixed-effects, and proxy-variable estimators reported in the next section. The Spearman correlation structure (Appendix Figure~\ref{fig:spearman_correlation}) preserves the strong positive associations among revenue, cost of sales, and labour costs that are characteristic of a coherent production technology, and the weak correlation of R\&D with the standard inputs is consistent with innovation being a margin distinct from input scale.

Significantly, the R\&D variable retains its characteristic zero-dominated shape: only $2.0\%$ of firm-years carry a strictly positive R\&D entry, the median R\&D stock is exactly zero, and the right tail of $\ln R$ accounts for the entirety of the within-panel variation in innovation intensity (skewness of $7.75$ in the log stock). This structural feature is the empirical reason for the $\log(R+1)$ transformation rather than a plain logarithm, and it motivates treating R\&D as an endogenous control over the productivity transition function in the \textcite{doraszelski2013r} specification rather than a continuous input on par with $K$, $L$, and $M$.

Three dimensions of heterogeneity in the estimation sample merit particular attention. First, at the firm level (Table~\ref{tab:desc_log}), the dispersion in capital intensity ($\sigma_{\ln K} = 2.626$) substantially exceeds that of output ($\sigma_{\ln Y} = 2.042$), reflecting marked differences in production technology across firms of similar revenue scale. Second, at the sectoral level (Appendix Figure~\ref{fig:rd_share_by_sector}), R\&D engagement is concentrated in information and communication and manufacturing, whereas extractive industries, agriculture, and utilities exhibit near-zero innovation rates; this sectoral polarization is the empirical hook for the cross-sectoral analysis reported in Section: Cross-Sectoral Analysis and for the heterogeneous-returns interpretation of the \textcite{doraszelski2013r} marginal-effect curves. Third, on a temporal level (Appendix Figure~\ref{fig:box_revenue_rd}), the panel covers a window of significant macroeconomic instability, ranging from the 2014-2015 ruble devaluation and the consecutive sanction regimes to the post-2022 geopolitical shift; this period offers the exogenous environmental variation necessary to identify productivity dynamics across the 2011-2023 horizon.


\section{Results}

This section presents the empirical findings of the study. The discussion proceeds from the main production function estimates and their economic interpretation, through a graphical analysis of productivity dynamics, to a brief assessment of the diagnostic tests.

\subsection{Production Function Estimates}

\begin{table}[htbp]
\centering
\caption{Production Function Estimates and Cross-Model Diagnostics}
\label{tab:main_results}
\begin{tabular}{lcccccc}
\toprule
 & OLS & FE & FE+time & OP & LP & DJ \\
\midrule
\multicolumn{7}{l}{\textit{Input elasticities}} \\
  Labor ($\ln L$) & $0.0892^{***}$ & $0.1365^{***}$ & $0.1376^{***}$ & $0.0690^{***}$ & $0.1307^{***}$ & $0.1306^{***}$ \\
   & $(0.0015)$ & $(0.0021)$ & $(0.0023)$ & $(0.0039)$ & $(0.0094)$ & $(0.0076)$ \\
  Materials ($\ln M$) & $0.8805^{***}$ & $0.8030^{***}$ & $0.8025^{***}$ & $0.8827^{***}$ & $0.8414^{***}$ & $0.8420^{***}$ \\
   & $(0.0018)$ & $(0.0029)$ & $(0.0029)$ & $(0.0048)$ & $(0.0090)$ & $(0.0074)$ \\
  Capital ($\ln K$) & $0.0044^{***}$ & $0.0260^{***}$ & $0.0251^{***}$ & $0.0493^{***}$ & $0.0010^{***}$ & $0.0010^{***}$ \\
   & $(0.0006)$ & $(0.0009)$ & $(0.0009)$ & $(0.0056)$ & $(0.0002)$ & $(0.0001)$ \\
  R\&D Stock ($\ln R$) & $0.0209^{***}$ & $0.0037^{*}$ & $0.0046^{**}$ & $0.0123^{**}$ & $0.0150^{***}$ & --- \\
   & $(0.0015)$ & $(0.0019)$ & $(0.0019)$ & $(0.0099)$ & $(0.0018)$ &  \\
\midrule
\multicolumn{7}{l}{\textit{Technology and TFP dynamics}} \\
  Returns to Scale ($\sum \beta$) & 0.974 & 0.966 & 0.965 & 1.001 & 0.973 & 0.974 \\
  \quad Classification & CRS & DRS & DRS & CRS & CRS & CRS \\
  Persistence ($\rho$) & 0.7193 & 0.1167 & 0.1161 & 0.9872 & 0.7302 & 0.7303 \\
  $\mathbb{E}[\xi_{it}]$ & --- & --- & --- & -0.0000 & -0.0000 & -0.0000 \\
\midrule
\multicolumn{7}{l}{\textit{Fit and GMM diagnostics}} \\
  Observations & 691,644 & 691,644 & 691,644 & 64,453 & 691,644 & 691,644 \\
  $R^2$ & 0.9359 & 0.8349 & 0.8302 & --- & --- & --- \\
  Hansen $J$ ($p$-value) & --- & --- & --- & --- & 4.35 & 8.14 \\
    &  &  &  &  & $p$=0.114 & $p$=0.017 \\
  $J/n$ & --- & --- & --- & --- & 0.000008 & 0.000015 \\
\bottomrule
\end{tabular}
\begin{flushleft}
\footnotesize
\textit{Note:} Clustered bootstrap standard errors in parentheses (firm-level resampling).
$^{*}$ $p<0.1$; $^{**}$ $p<0.05$; $^{***}$ $p<0.01$.
OLS uses HC3 robust standard errors.
FE = Entity fixed effects; FE+time = Two-way (entity + time) fixed effects.
OP = Olley--Pakes; LP = Levinsohn--Petrin (materials as proxy; $\beta_m$ identified via ACF; R\&D treated as a state variable so $\beta_r$ is identified jointly with $\beta_k$);
DJ = Doraszelski--Jaumandreu (R\&D enters the productivity transition function rather than the production function).
RTS = $\sum_{j \in \{l, m, k\}} \beta_j$ (R\&D excluded from RTS, following standard practice).
$\rho$ = AR(1) coefficient of the model-specific $\omega_{it}$.
$\mathbb{E}[\xi_{it}]$ = sample mean of the TFP innovation (zero by construction in moment-based estimators).
\end{flushleft}
\end{table}

Table~\ref{tab:main_results} \hypertarget{target:main_results}{} provides a comprehensive overview of the production function estimation results across six specifications: Ordinary Least Squares (OLS), entity Fixed Effects (FE), two-way Fixed Effects with time dummies (FE+time), \textcite{olley1992dynamics} (OP), \textcite{levinsohn2003estimating} (LP), and \textcite{doraszelski2013r} (DJ). All coefficients on the standard production inputs are statistically significant at the 1\% level, confirming that labour, capital, and materials are meaningful determinants of firm output across all model specifications.

The labour elasticity ranges from 0.069 (OP) to 0.138 (FE+time), with OLS at 0.089. The upward move from OLS/OP to FE is consistent with a downward simultaneity bias in the pooled estimates once unobserved firm heterogeneity is absorbed; the LP and DJ estimates coincide at 0.131, consistent with both estimators pulling labour upward through the proxy-variable correction.

Capital elasticities vary more sharply. OLS (0.004) and the FE estimators (${\sim}0.025$) likely understate $\beta_k$ through the classic transmission bias. However, the OP estimator recovers a considerably larger coefficient (0.049), while LP and DJ return near-zero values (0.001 in both cases); this compression is a known feature of control-function estimators when the proxy (materials under LP, the $\phi$-inversion augmented with lagged R\&D under DJ) absorbs much of the variation attributable to capital.

Materials ($\ln M$) consistently come out as the dominant input, with elasticities ranging from 0.8025 (FE+time) to 0.8827 (OP). This finding is consistent with the cost structure of Russian industrial firms, where intermediate consumption forms the largest share of total production costs. The stability of the materials coefficient across the OLS, FE, \textcite{olley1992dynamics}, and \textcite{doraszelski2013r} specifications (approximately 0.80 to 0.88) lends credibility to the underlying production technology and indicates that the Cobb-Douglas functional form provides a reasonable approximation for this input.

Of particular interest for the research questions posed in this study is the coefficient on the R\&D knowledge stock ($\ln R$). In the OLS specification, R\&D enters with a positive and highly significant elasticity of 0.0209 ($p < 0.01$), implying that a 1\% increase in the accumulated R\&D stock is associated with approximately a 0.021\% increase in firm revenue, ceteris paribus. Although this magnitude appears humble, it is broadly consistent with the range of 0.07 to 0.15 reported in the international literature for developed economies \parencite{Pieri2018Modelling, Kancs2016R&D}, adjusting for the fact that the Russian sample is dominated by non-innovators; their zero R\&D stock artificially deflates the average elasticity. Upon introducing entity fixed effects, the R\&D coefficient declines to $0.0037$ ($p < 0.1$), and with two-way fixed effects it edges to $0.0046$ ($p < 0.05$); both estimates are positive and at the lower edge of statistical significance, an order of magnitude below the pooled OLS coefficient. This attenuation is notable: it suggests that a substantial portion of the cross-sectional association between R\&D and output reflects persistent, time-invariant firm characteristics, such as managerial quality or strategic orientation, rather than a within-firm causal effect of knowledge accumulation. The \textcite{olley1992dynamics} structural estimate of $0.0123$ remains positive ($p < 0.05$) but is estimated with a comparatively large standard error ($0.0099$), so the within-firm causal interpretation has to be tempered by the limited precision of the OP point estimate. This imprecision potentially arises from the substantial loss of observations necessitated by the OP framework, which reduces the estimation sample from $691{,}644$ firm-years to a markedly smaller subset owing to the requirement for strictly positive investment. The \textcite{levinsohn2003estimating} estimator, which retains all firms regardless of investment status, yields $\beta_r = 0.0150$ ($p < 0.01$) on the full sample of $691{,}644$ firm-years; the corresponding Hansen $J$ test does not reject the over-identifying restrictions ($p = 0.114$), supporting the validity of the moment conditions used to identify $\beta_r$.

In the \textcite{levinsohn2003estimating} model the R\&D capital stock $r_{it}$ enters the production function on the same footing as physical capital, it is a state variable evaluated at $t-1$, not a contemporaneous response to the productivity shock, and its coefficient $\beta_r$ is identified jointly with $(\beta_l, \beta_m, \beta_k)$ from the Stage~2 moment conditions, exactly as $\beta_k$ is. The role of materials in the LP specification is to invert out unobserved productivity $\omega_{it}$ via the proxy polynomial $\phi_t(l, m, k, r)$, not to crowd R\&D out of the production function: conditioning $\phi_t$ on $r_{it}$ ensures that the recovered $\omega_{it}$ is purged of the R\&D-driven component of the firm's decision rule, so the lagged moment conditions used to identify $\beta_r$ are valid (Section Models). Under the DJ framework, by contrast, R\&D does not enter the production function at all; it is modelled as a forward-looking choice that shifts the distribution of future productivity through the controlled Markov law $\omega_{it} = g(\omega_{it-1}, r_{it-1}) + \xi_{it}$, so including $r_{it}$ as a contemporaneous input alongside $g(\cdot)$ would double-count its effect and violate the scalar-unobservable restriction required for identification. The two specifications are therefore complementary: LP identifies the static output elasticity of R\&D capital under proxy-variable corrections, while DJ identifies the dynamic transition of productivity in response to lagged R\&D.

For this research the specification is also substantively appropriate. The empirical setting is a large, heterogeneous panel of Russian industrial firms in which reported R\&D is highly uneven, heavily under-reported in conventional accounting lines, and dominated by zeros (non-innovators). These features are exactly those for which the reduced-form elasticity on $\ln R$ in OLS and FE is known to be attenuated and difficult to interpret causally. Routing R\&D through the transition function $g(\cdot)$, as in DJ, accommodates the zero mass naturally, recovers the full dynamic payoff of innovation (rather than a static elasticity contaminated by simultaneity), and is consistent with the way innovation is financed and deployed in practice: R\&D today modifies the firm's future productivity draw, not today's output given fixed inputs. The LP specification, in turn, provides a clean benchmark in which materials absorb the contemporaneous informational content of productivity without imposing a structural role for R\&D, which is the correct ``agnostic'' comparator for the DJ model. Overall, the two specifications are therefore complementary rather than rival, and both are appropriate for the data at hand.

Although the reduced-form models estimate the contemporaneous output elasticity of knowledge capital, the \textcite{doraszelski2013r} framework captures the dynamic, forward-looking effect of R\&D on the evolution of total factor productivity over time.

Importantly, the returns to scale (RTS), computed as the sum of input elasticities, provide further insight into the production technology. The OLS, FE, and \textcite{doraszelski2013r} models yield RTS values close to but slightly below unity (0.965 to 0.974), suggesting mildly decreasing returns to scale. The \textcite{olley1992dynamics} specification returns an RTS of 1.001, indicating effectively constant returns. Whereas the \textcite{levinsohn2003estimating} estimate of 0.975 is consistent with the remaining specifications and confirms slight decreasing returns under the ACF identification strategy.

\paragraph{Robustness and diagnostics.} The specification passes the diagnostic checks one would expect of a structural production-function model at this sample size: the OLS baseline achieves $R^2 = 0.936$ with all VIFs well below $10$, classical heteroscedasticity and serial-correlation tests are rejected (as is typical in very large panels) and are addressed through HC3 and clustered bootstrap standard errors, and the Hansen $J$-tests for LP and DJ reject formally but yield scaled statistics $J/n<10^{-4}$, consistent with negligible per-observation misspecification. Three further robustness checks confirm that the headline estimates are not artefacts of specific instrument or polynomial choices. The LP lag comparison (Appendix~\ref{app:diagnostics}, Table~\ref{tab:lp_lag}) shows that moving from the ACF-default one-lag instrument set to a two-lag set leaves the input elasticities and the R\&D coefficient within approximately one bootstrap standard error. The polynomial-sensitivity exercise for OP and LP (Appendix~\ref{app:diagnostics}, Table~\ref{tab:poly_sensitivity}) shows that the first-stage control function is stable across polynomial orders two to four in the proxy variables. The DJ polynomial-order sensitivity check (Appendix~\ref{app:dj_polynomial_order}) confirms that the input elasticities and returns to scale are numerically close under the parsimonious degree-2 specification and the preferred degree-3 specification, so the case for the richer specification rests on the heterogeneity of the cross-derivative that only the degree-3 polynomial can identify rather than on point-estimate differences. Full residual plots, VIF tables, LP lag comparisons, and polynomial-sensitivity results are reported in Appendix~\ref{app:diagnostics}.


\paragraph{The Dynamic Returns to R\&D: \textcite{doraszelski2013r} Transition Function}

\begin{table}[htbp]
\centering
\caption{DJ Transition Function (Degree-3 Polynomial): $E[\omega_t \mid \omega_{t-1}, r_{t-1}]$}
\label{tab:dj_transition}
\begin{tabular}{lr}
\toprule
Parameter & Estimate \\
\midrule
  $\gamma_0$ (constant) & $0.0098$ \\
  & $\text{(0.0220)}$ \\
  $\gamma_1$ ($\omega_{t-1}$) & $1.1897^{***}$ \\
  & $\text{(0.0549)}$ \\
  $\gamma_2$ ($r_{t-1}$) & $-0.0004$ \\
  & $\text{(0.0211)}$ \\
  $\gamma_3$ ($\omega_{t-1}^2$) & $-0.2809^{***}$ \\
  & $\text{(0.0345)}$ \\
  $\gamma_4$ ($r_{t-1}^2$) & $0.0014$ \\
  & $\text{(0.0036)}$ \\
  $\gamma_5$ ($\omega_{t-1} \times r_{t-1}$) & $0.0080$ \\
  & $\text{(0.0269)}$ \\
  $\gamma_6$ ($\omega_{t-1}^3$) & $0.0299^{***}$ \\
  & $\text{(0.0050)}$ \\
  $\gamma_7$ ($r_{t-1}^3$) & $-0.0001$ \\
  & $\text{(0.0001)}$ \\
  $\gamma_8$ ($\omega_{t-1}^2 \times r_{t-1}$) & $-0.0124^{*}$ \\
  & $\text{(0.0058)}$ \\
  $\gamma_9$ ($\omega_{t-1} \times r_{t-1}^2$) & $0.0003$ \\
  & $\text{(0.0043)}$ \\
\midrule
GMM criterion & 0.000015 \\
Hansen $J$ & 8.1413 \\
\quad d.f. & 2 \\
\quad $p$-value & 0.0171 \\
\bottomrule
\end{tabular}
\begin{flushleft}
\footnotesize
\textit{Note:} Degree-3 complete polynomial in $(\omega_{t-1}, r_{t-1})$ with 10 parameters,
following \textcite{doraszelski2013r}.
Two-step GMM estimates with profiled $\gamma$ (ACF 2015).
Bootstrap standard errors in parentheses.
$^{*}$ $p<0.1$; $^{**}$ $p<0.05$; $^{***}$ $p<0.01$.
\end{flushleft}
\end{table}

The structural contribution of this research lies in the estimated productivity transition function $\omega_{it} = g(\omega_{it-1}, r_{it-1}) + \xi_{it}$, specified as a complete third-degree polynomial in $(\omega_{it-1}, r_{it-1})$ with ten parameters following the original \textcite{doraszelski2013r} formulation. Table~\ref{tab:dj_transition} \hypertarget{target:dj_transition}{} reports the parameter estimates; the corresponding sensitivity check against a parsimonious degree-2 specification is reported in Appendix~\ref{app:dj_polynomial_order} and confirms that the degree-3 polynomial is needed only to identify the heterogeneity of the cross-derivative discussed below.

Three substantive features emerge from Table~\ref{tab:dj_transition}. First, the linear persistence coefficient $\gamma_1$ is large and positive, indicating that productivity is highly persistent at the firm level: most of the current $\omega$ is carried forward into the next period through the AR component of the transition. The non-linear persistence terms $\gamma_3$ and $\gamma_6$ are negative and statistically significant, generating a concave reversion profile in which firms at the upper tail of the productivity distribution mean-revert more strongly than those at the median, a quantitatively meaningful refinement of the linear AR benchmark and a feature that a quadratic polynomial cannot accommodate. Second, the direct R\&D coefficient $\gamma_2$ is small and statistically indistinguishable from zero in the bootstrap, but the higher-order R\&D terms $\gamma_4$ and $\gamma_7$ together with the interaction coefficients $\gamma_5$, $\gamma_8$ and $\gamma_9$ are jointly informative: the marginal effect of lagged R\&D on expected productivity is recovered not from any single coefficient but from the full polynomial gradient evaluated at the within-sample distribution of $(\omega, r)$. Third, the negative sign of $\gamma_8$ together with the positive cubic persistence term $\gamma_6$ produces a transition surface in which R\&D and current productivity are state-dependent substitutes throughout the in-sample support, with the strength of substitution growing as the firm moves towards the productivity frontier.

The policy-relevant marginal effect $\partial \mathbb{E}[\omega_t]/\partial r_{t-1} = \gamma_2 + 2\gamma_4 r + \gamma_5 \omega + 3\gamma_7 r^2 + 2\gamma_9 \omega r + \gamma_8 \omega^2$, evaluated at the mean log R\&D among innovators ($\bar{r} = 3.93$) and at sample quantiles of $\omega_{it-1}$, is tabulated in Appendix~\ref{app:diagnostics} (Table~\ref{tab:dj_marginal} \hypertarget{target:dj_marginal}{}). The marginal return to lagged R\&D is positive and economically meaningful for firms at the lower deciles of the productivity distribution and decays monotonically as $\omega$ increases, declining from $+0.008$ at the 10th percentile to $+0.001$ at the 90th percentile. The cross-derivative $\partial^2 g / \partial\omega\partial r = \gamma_5 + 2\gamma_8\omega + 2\gamma_9 r$ is the central object of the structural test: it is uniformly negative on the in-sample support, and its absolute magnitude grows monotonically across $\omega$-quantiles (from $|-0.007|$ in the bottom decile to $|-0.021|$ in the top decile), implying that R\&D and current productivity are increasingly strong substitutes as one moves towards the technological frontier (Appendix~\ref{app:diagnostics}, Tables~\ref{tab:dj_cross_deriv} and~\ref{tab:dj_cd_by_quantile}).

\begin{figure}[htbp]
   \centering
   \includegraphics[width=0.78\textwidth]{dj_marginal_effect_rd_degree3.png}
   \caption{Marginal effect of lagged R\&D on expected productivity, $\partial \mathbb{E}[\omega_{it}]/\partial r_{it-1}$, from the degree-3 \textcite{doraszelski2013r} transition function. The three curves correspond to the 10th, 50th, and 90th percentiles of the within-sample productivity distribution; the shaded band is the 95\% bootstrap confidence interval at the median $\omega$, and the vertical dotted line marks the median log R\&D stock among innovators. Curves fan out across $\omega$-quantiles - a heterogeneity in the marginal effect that the quadratic specification cannot identify by construction.}
   \label{fig:dj_marginal_effect_rd}
\end{figure}

Figure~\ref{fig:dj_marginal_effect_rd} \hypertarget{target:dj_marginal_effect}{} renders the same finding graphically. The three curves are non-parallel and the gap between them widens with the lagged R\&D stock, both features being a direct empirical consequence of the cubic interaction terms $\gamma_7 r^2$, $\gamma_8 \omega^2$ and $\gamma_9 \omega r$. Economically, the picture is one of state-dependent absorptive capacity: firms operating below the sectoral productivity median capture a positive marginal payoff from each additional unit of R\&D capital, the payoff attenuates monotonically along the $\omega$-axis, and the upper-decile firms are at the inflection beyond which lagged R\&D ceases to add to expected productivity. This pattern, recovered from the structural polynomial without imposing it ex ante, is consistent with the distance-to-frontier prediction of step-by-step Schumpeterian competition \parencite{aghion1990model}. It is also worth noting that the sign of the cross-derivative is reversed relative to the original \textcite{doraszelski2013r} estimates for Spanish manufacturing, where lagged productivity and R\&D were found to be predominantly complements; the substantive interpretation of this contrast is reserved for the corresponding paragraph of the Discussion section.

\begin{figure}[htbp]
   \centering
   \includegraphics[width=0.8\textwidth]{rd_premium_DJ.png}
   \caption{The R\&D productivity premium estimated using the DJ framework, highlighting the efficiency gap between innovators and non-innovators over the period 2012-2023.}
   \label{fig:rd_premium_dj}
\end{figure}

The patterns observed in the figures mirror and extend the insights gleaned from the tabular analysis, ensuring a robust verification of the results. Figure~\ref{fig:rd_premium_dj} \hypertarget{target:fig1}{} presents a yearly comparison of TFP scores between innovative firms (those with positive R\&D expenditure) and conventional firms, estimated using the \textcite{doraszelski2013r} model over the period 2012 to 2023. Throughout the entire observation window, firms engaging in R\&D consistently exhibit higher median TFP than their non-innovating counterparts. Indeed, this "R\&D premium" holds firm even in the early years of the panel and appears to widen modestly over time, suggesting that the cumulative nature of knowledge capital supports productivity differentials. Although this descriptive pattern does not, in itself, establish causality, it provides corroborative evidence consistent with the structural estimates reported in Table~\ref{tab:dj_transition}.

\begin{figure}[htbp]
   \centering
   \includegraphics[width=0.8\textwidth]{tfp_distribution.png}
   \caption{Distribution of Total Factor Productivity (TFP) scores estimated via various econometric models, illustrating the density of productivity across OLS, FE, and proxy-based methods.}
   \label{fig:tfp_distribution}
\end{figure}

Figure~\ref{fig:tfp_distribution} \hypertarget{target:fig3}{} depicts the comparative distribution of cumulative TFP estimates across all five models (OLS, FE, OP, LP, DJ). The density plots reveal that the OLS and FE distributions are relatively concentrated, whereas the structural estimators, particularly \textcite{levinsohn2003estimating} and \textcite{doraszelski2013r}, produce wider dispersions. Such variability underscores the structural models' ability of the structural models to recover a broader range of firm-level productivity heterogeneity that is otherwise obscured by the simultaneity bias inherent in naive estimators. The \textcite{doraszelski2013r} distribution, in particular, displays a rightward shift relative to the FE estimates, indicating that once the endogenous accumulation of knowledge capital is accounted for, the recovered productivity distribution more accurately captures the efficiency gains attributable to innovation.

\begin{figure}[htbp]
   \centering
   \includegraphics[width=0.8\textwidth]{tfp_evolution.png}
   \caption{Dynamics of median Total Factor Productivity (TFP) of Russian industrial firms (2011-2023), comparing the evolution of productivity estimates across different model specifications.}
   \label{fig:tfp_evolution}
\end{figure}

Figure~\ref{fig:tfp_evolution} \hypertarget{target:fig4}{} traces the time-series evolution of median log-TFP for Russian industrial firms from 2012 to 2023, highlighting the sensitivity of the recovered productivity series to the choice of estimation methodology. The six estimators arrange themselves into three clearly separated bands. At the top of the figure ($\omega \approx 1.02$, almost completely flat), the entity fixed-effects (FE) and two-way fixed-effects (FE+time) trajectories sit on top of one another, visually a single line, reflecting the fact that within-firm demeaning absorbs most of the temporal variation in productivity into the firm-specific intercept rather than into a year-by-year median path. In the middle band ($\omega \approx 0.62$ in 2012, drifting upwards), the pooled OLS series describes a smooth, monotone climb to $\omega \approx 0.73$ by 2023 - a trajectory that, in the absence of any correction for input--productivity simultaneity, conflates input expansion with genuine technical change and is therefore the natural upper bound of the productivity-growth narrative one would tell from the raw data. The \textcite{olley1992dynamics} (OP) series, by contrast, sits about one decimal below OLS at $\omega \approx 0.62$--$0.64$ and is essentially flat across the panel; this reflects a combination of the OP positive-investment sample restriction and the absorption of the trend component into the proxy-controlled $\phi$ inversion. In the bottom band ($\omega \approx -0.07$ in 2012, rising to $\omega \approx +0.04$ by 2023), the proxy-variable Levinsohn--Petrin (LP) and the structural \textcite{doraszelski2013r} (DJ) series are visually indistinguishable on the median: both estimators recover a productivity distribution that is centred close to zero and that displays a slightly negative drift through 2020 followed by a modest post-pandemic rebound. The vertical separation between the bottom band (LP, DJ), the middle band (OLS, OP), and the top band (FE, FE+time) is the headline empirical signature of identification: estimators that absorb input--productivity simultaneity assign more of the observed output dispersion to inputs (and therefore less to $\omega$), whereas the FE estimators that absorb only firm-level heterogeneity place most of the residual output variance on the productivity series itself. The descriptive widening of the LP/DJ series after 2014 and again after 2020 is suggestive of the catch-up dynamics documented in the cross-sectoral analysis but does not, on its own, identify a causal sanctions effect; the structural interpretation of the post-2014 R\&D premium is reserved for the corresponding paragraph of the Discussion section.

\subsection{Cross-sectoral Analysis}
\label{sec:cross_sectoral}

Estimating each of the six specifications (OLS, entity FE, entity-and-time FE, OP, LP, and DJ) separately on every industry produces a coherent cross-sectoral picture. In every method materials remain the dominant input in every sector, with elasticities concentrated between 0.55 and 0.97 once the proxy-variable correction is in place; labour shares are largest in information and communication and transportation, while capital coefficients reach their highest values under OP, consistent with the textbook simultaneity correction for forward-looking input choice. The within-firm R\&D coefficient turns positive and economically meaningful in mining, transportation, and manufacturing once unobserved heterogeneity is absorbed, while it collapses to zero or negative values in agriculture, construction, electric energy and gas, water supply, and information and communication are the five sectors in which the cross-sectional R\&D-output association is mostly absorbed by time-invariant firm characteristics. The full per-sector tables and coefficient figures for each of the six methods are reported in Appendix~\ref{app:industry_methods}.\hypertarget{target:industry_coefficients_ols}{}\hypertarget{target:industry_coefficients_fe}{}\hypertarget{target:industry_coefficients_fe_tw}{}\hypertarget{target:industry_coefficients_op}{}\hypertarget{target:industry_coefficients_lp}{}\hypertarget{target:industry_coefficients_dj}{}

\begin{figure}[htbp]
   \centering
   \includegraphics[width=0.8\textwidth]{industry_rts_comparison.png}
   \caption{Comparison of Returns to Scale (RTS) across industries, illustrating the sum of labour and capital elasticities by estimation method.}
   \label{fig:industry_rts_comparison}
\end{figure}

Figure~\ref{fig:industry_rts_comparison} \hypertarget{target:industry_rts_comparison}{} provides a comparative visualization of returns to scale across the eight production sectors and six estimation methods. Three patterns stand out. (a)~The \emph{rank ordering} of sectors by RTS is remarkably stable across estimators: agriculture and mining consistently sit at the upper end (RTS in 0.93--1.06 across methods), while transportation and information and communication occupy the lower end (0.83--0.98), reflecting the difficulty of scaling spatially dispersed logistics networks and knowledge-intensive services respectively. (b)~The \emph{level} of RTS is systematically lower under the proxy-variable estimators (LP and DJ), which is exactly the direction predicted by the correction for input--productivity simultaneity: absorbing the anticipated-shock component of variable inputs into $\omega$ mechanically compresses the input-elasticity sum. (c)~The OP estimator produces the most compressed RTS range across sectors because of the investment-positivity sample restriction, a known cost of the OP approach at industry granularity.

The industry-level TFP distributions (Appendix~Figures~\ref{fig:industry_tfp_ols} \hypertarget{target:industry_tfp_ols}{},~\ref{fig:industry_tfp_distribution_lp} \hypertarget{target:industry_tfp_distribution_lp}{},~\ref{fig:industry_tfp_distribution_dj} \hypertarget{target:industry_tfp_distribution_dj}{},~\ref{fig:industry_tfp_op} \hypertarget{target:industry_tfp_op}{}) reveal substantial cross-sectoral heterogeneity in productivity levels. Under the structural estimators (LP and DJ), the TFP distributions of manufacturing and construction are relatively concentrated, while mining and agriculture exhibit markedly wider dispersions, indicating greater within-sector heterogeneity in production efficiency, a pattern consistent with the commodity-price exposure of extractive industries and the large productivity gap between mechanised and non-mechanised producers in Russian agriculture. The median TFP ranks mining and information and communication at the top and water supply at the bottom across estimators, with the DJ distributions additionally compressed relative to OLS (a smaller right tail and a tighter interquartile range), confirming that routing R\&D through the productivity transition function re-attributes right-tail dispersion from unexplained residuals to the cumulative effect of innovation.

Taken together, the cross-sectoral picture under the structural estimators (LP and DJ) yields four principal conclusions. The reduced-form OLS and FE results, which are biased by input--productivity simultaneity, are reported sector-by-sector in Appendix~\ref{app:industry_methods} but are not used as a basis for the substantive interpretation that follows.
\begin{enumerate}
    \item \textbf{An IT paradigm shift: from raw materials to human capital.} Although the reduced-form FE and OLS estimators wash out cross-sectoral differences, the structural LP and DJ specifications agree strikingly on the technology of the information and communication sector. Whereas in classical manufacturing the materials share rises to $\beta_m \approx 0.87$, in information and communication it falls to $\beta_m \approx 0.66$ and the labour elasticity climbs to $\beta_l = 0.27$ under LP and $\beta_l = 0.31$ under DJ, among the highest values recorded across the eight sectors. This sectoral signature is the empirical fingerprint of a structural transformation: in the digital economy the marginal value-added is created by human capital (engineers, developers, designers) rather than by intermediate physical inputs, in sharp contrast with the materials-intensive technology of traditional industry.

    \item \textbf{The R\&D returns paradox: infrastructure outperforms high-tech.} Reading the LP specification (which, unlike DJ, identifies a static R\&D elasticity) yields a counter-intuitive sectoral ranking. The largest output elasticity of the R\&D capital stock is observed not in manufacturing ($\beta_r = 0.003$) or in information and communication ($\beta_r = 0.005$), but in transportation ($\beta_r = 0.041$) and water supply ($\beta_r = 0.022$). The economic reading is that mature production-and-tech sectors are operating in the diminishing-marginal-returns region of the innovation curve, each additional unit of R\&D capital is incremental and yields only a small productivity gain, while sluggish infrastructure sectors (transport, utilities) start from such a low technological baseline that the introduction of even basic R\&D-enabled solutions, such as routing algorithms in transport or smart leak-detection in water supply, generates an outsize productivity payoff. The interpretation is consistent with the broader catch-up logic underpinning the absorptive-capacity finding from the DJ transition function.

    \item \textbf{The ``ghost of capital'': extreme rigidity of the fixed-asset stock.} Both proxy-variable specifications, which correct for the input--productivity simultaneity that biases naive estimators, expose a striking macro-level (or behavioural) feature of the Russian production economy: physical capital exerts almost no contemporaneous effect on output, with $\beta_k$ estimates concentrated in the $0.001$--$0.03$ range across all eight sectors. Where pooled OLS and FE attempt to assign capital at least some role, the structural correction reveals that within-firm output variation is driven essentially by intermediate inputs ($\beta_m$) and, to a lesser extent, by labour. Three economic interpretations are mutually compatible: substantial under-utilisation of large capital stocks accumulated in earlier investment cycles, deep depreciation of the existing fixed-asset base unaccounted for in book values, and a binding inelasticity of capital-investment budgets at the firm level - so that growth, when it occurs, is financed through working capital rather than new fixed investment.

    \item \textbf{Decreasing returns to scale in logistics as a marker of market frictions.} The cross-sectoral RTS comparison delivers a sharp and unambiguous result for transportation: under both LP and DJ, returns to scale fall to $\text{RTS} = 0.825$ - the most pronounced decreasing-returns finding in the panel, against the near-constant returns recorded in the remaining seven sectors. In applied industrial-organisation work, RTS substantially below unity is read not as a pure technological constraint but as an indicator of binding market frictions or capacity constraints: as transport firms scale, they encounter rising transaction and managerial costs (congested networks, scarce warehousing capacity, growing logistical complexity) that erode the marginal productivity of every additional rouble of input. The combination of a high $\beta_r$ and a low RTS in transportation thus paints a coherent economic picture: large potential payoffs to R\&D are available, but they are absorbed by the very frictions that R\&D would help relax.
\end{enumerate}

\section{Discussion}

The headline empirical result of this study is twofold. First, on the pooled panel of $691{,}644$ firm-year observations covering eight production-oriented OKVED sections of the Russian economy, structural identification of the R\&D capital elasticity yields $\beta_r^{\text{LP}} = 0.0150$ ($p<0.01$) and $\beta_r^{\text{OP}} = 0.0123$ ($p<0.01$); both are positive, statistically significant, and located within the bootstrap confidence band that excludes zero. Second, the \textcite{doraszelski2013r} transition function reveals that this average effect masks a substantial degree of heterogeneity: the marginal effect of lagged R\&D on expected productivity is largest for firms in the lower deciles of the within-sample $\omega$-distribution and decays monotonically as the firm approaches the sectoral frontier. The remainder of this section interprets these two findings through the theoretical lenses adopted in Section~\ref{sec:literature}, places them alongside prior empirical evidence on Russian and international firms, and draws out the implications for innovation policy and productivity forecasting.

\paragraph{State-dependent returns and the sign reversal relative to Spanish manufacturing.} The empirical pattern recovered from the degree-3 transition function - positive, decreasing marginal returns to R\&D as firm productivity rises towards the sectoral frontier - is the central novelty of the study and has a direct theoretical interpretation. It maps onto the distance-to-frontier prediction of step-by-step Schumpeterian competition \parencite{aghion1990model}: firms that lag the technological frontier benefit disproportionately from each marginal unit of R\&D capital because the catch-up margin is wide and the knowledge they acquire substitutes for an in-house knowledge gap. Conversely, frontier-adjacent firms exhaust the catch-up channel and additional R\&D outlays compete with capital and human-resource alternatives that may carry higher marginal returns at the frontier. Structurally, this manifests in estimates as a uniformly negative cross-derivative $\partial^2 g/\partial\omega\partial r$ whose absolute magnitude grows monotonically across $\omega$-quantiles (Appendix~\ref{app:diagnostics}, Table~\ref{tab:dj_cd_by_quantile}).

The sign of this cross-derivative is the opposite of what \textcite{doraszelski2013r} report for Spanish manufacturing, where lagged productivity and R\&D act predominantly as \emph{complements} (in seven of their nine industries). The reversal is an empirical finding rather than an artefact: the Spanish panel is concentrated in manufacturing firms with established in-house R\&D programmes, a setting in which prior productivity reinforces the marginal payoff to R\&D; the Russian sample is dominated by non-innovators, spans extractive, utility, transportation, and selected service activities alongside manufacturing, and covers a post-sanctions regime in which catch-up to a moving foreign frontier is the operative margin. Under these conditions the marginal payoff to lagged R\&D is dominated by the catch-up channel rather than by the reinforcement channel, and the cross-derivative inverts.

\paragraph{Comparison with prior empirical evidence.} The pooled-panel R\&D elasticities recovered here ($\beta_r^{\text{OLS}} \approx 0.021$, $\beta_r^{\text{LP}} \approx 0.015$, $\beta_r^{\text{OP}} \approx 0.012$) sit below the range of $0.07$--$0.15$ commonly reported for developed-economy panels \parencite{Pieri2018Modelling, Kancs2016R&D}. Two features of the estimation environment account for the gap. First, only $2.0\%$ of firm-years carry a strictly positive R\&D entry on line 1120 of the statutory balance sheet (Section~Variables Construction); the average elasticity is therefore mechanically attenuated by the large mass of non-innovators whose zero R\&D stock pulls the linear coefficient towards zero, even though the structural transition function (the DJ specification) shows that conditional on positive innovation the marginal effect is comparable to the international benchmarks at the lower productivity quantiles. Second, line 1120 itself is known to under-record R\&D activity because firms exercise discretion in capitalising R\&D as intangible assets or expensing it through the income statement; on richer innovation proxies (patents, R\&D-survey variables) the elasticity in Russian samples typically rises towards the international range. Within the small Russian-microdata literature on R\&D and TFP, our estimates are broadly consistent with \textcite{malein2019total}, who report a $24$--$28\%$ TFP gain on adoption of an electric-arc-furnace technology in a 39-firm panel of Russian metallurgy, a within-firm technology shock whose magnitude is consonant with the upper-quantile catch-up implied by our DJ marginal-effect curve. Relative to the cross-sectional CDM-style work of \textcite{crepon1998research}, the DJ specification provides what those models cannot: an explicit, parametric statement about how R\&D shifts the \emph{distribution} of future productivity rather than the level alone.

\paragraph{R\&D premium dynamics across the sanctions regime.} Figure~\ref{fig:rd_premium_dj} traces the median TFP gap between innovators and non-innovators from 2012 to 2023 under the DJ specification. Throughout the observation period, the premium is consistently positive; it grew significantly after the 2014 sanctions episode and experienced a subsequent, albeit milder, expansion after 2022. Two readings of this temporal pattern bear directly on contemporary policy debates. On a Schumpeterian reading \parencite{aghion1990model}, sanctions raised the relative private return to in-house innovation by curtailing access to imported technology and to cross-border knowledge spillovers; firms that had been investing in R\&D before the shock retained their absorptive capacity and were better positioned to substitute domestic innovation for imported technology, thereby widening the productivity gap with non-innovators. This complementarity between in-house R\&D and external knowledge is precisely the channel documented in \textcite{Audretsch2020The} for UK firms, who find that R\&D and knowledge spillovers are jointly significant for labour productivity beyond their stand-alone effects; the post-2014 widening of the Russian R\&D premium is consistent with that mechanism operating in reverse, as the loss of cross-border spillovers raised the marginal value of the domestic R\&D component required to assimilate substitute technology. On a Griliches-style accounting reading \parencite{griliches1979issues}, the cumulative knowledge stock of innovating firms appreciated relative to that of their peers precisely because the depreciation of foreign-acquired technology accelerated under sanctions while the depreciation of home-grown knowledge did not. Although our data cannot causally separate these two channels, the widening of the premium is qualitatively consistent with both and provides empirical support for the import-substitution debate that has dominated Russian industrial-policy discussion since 2014.

\paragraph{Refining the theoretical lenses.} Two of the theoretical anchors adopted in Section~Literature Review need to be nuanced in light of the empirical evidence. First, the homogeneous-elasticity assumption of the original \textcite{griliches1979issues} knowledge-capital framework is rejected by the heterogeneity of $\partial \mathbb{E}[\omega_t]/\partial r_{t-1}$ across the in-sample distribution of $\omega$: the data require an interaction between productivity level and innovation effort that the standard log-linear Cobb--Douglas-with-knowledge-stock specification cannot accommodate. The empirical refinement we propose is to retain the Griliches accounting machinery (perpetual-inventory R\&D stock, $\beta_r$ as the static elasticity at the mean) but to read $\beta_r$ as an average over a non-uniform return profile rather than as a structural parameter that applies uniformly to all firms. Second, the Aghion--Howitt step-by-step competition framework \parencite{aghion1990model} predicts that R\&D returns are conditioned on the firm's distance to the technological frontier; our finding is exactly this prediction realised on Russian production-economy data, and to that extent the empirical evidence \emph{supports rather than contradicts} the Schumpeterian prior. The remaining theoretical work, which our empirical evidence calls for but does not itself supply, is to provide micro-foundations for the specific shape of the cross-derivative attenuation - in particular, why the marginal effect declines monotonically rather than U-shaped or inverted-U as some competition models would predict.

\paragraph{Implications for innovation policy and market efficiency.} Three implications for industrial policy and applied productivity work follow from the empirical results. First, blanket R\&D subsidies are dominated by targeted instruments: the marginal productivity payoff of an additional unit of R\&D support is maximised when the support is allocated to firms below the sectoral productivity median, where the absorptive-capacity margin is widest, and is smallest when it is allocated to incumbent frontier firms whose marginal effect has decayed towards zero. A uniform-rate R\&D credit thus systematically over-rewards inframarginal frontier innovation and under-rewards catch-up innovation, generating a static mis-allocation of innovation resources whose magnitude is bounded above by the spread of $\partial \mathbb{E}[\omega_t]/\partial r_{t-1}$ across $\omega$-quantiles documented in Figure~\ref{fig:dj_marginal_effect_rd}. Second, the persistence of the R\&D productivity premium under sanctions is consistent with a substitution of domestic innovation for restricted access to foreign technology and provides an empirical anchor for the import-substitution wing of contemporary Russian industrial policy. Third, for applied productivity forecasting and efficiency benchmarking, any model that imposes an exogenous first-order Markov law on $\omega_{it}$ will systematically overstate `passive' inertia and understate the returns to discretionary innovation effort, precisely because the share of persistence that DJ re-attributes to lagged R\&D becomes part of the AR coefficient under exogenous-Markov benchmarks. The corollary for forecasters is that out-of-sample TFP projections that ignore the R\&D channel will underestimate the productivity dispersion that endogenous innovation generates over a multi-year horizon.

\paragraph{Limitations.} Three limitations bound the scope of the inferences above. First, the pointwise $95\%$ bootstrap band on the marginal-effect curve is wide relative to its curvature, so the declining profile across $\omega$-quantiles should be read as suggestive rather than sharply identified at conventional confidence levels. Second, the analysis is confined to within-sample quantiles of the structural $\omega$, so the inferences describe heterogeneity present in the data rather than extrapolated frontier behaviour. Third, the R\&D measure inherits the line-1120 accounting discretion documented in the literature; on richer innovation proxies the magnitude of the elasticity is likely to rise, but the qualitative shape of the marginal-effect curve and the cross-sectoral ordering are unlikely to change.



\section{Conclusion}


The aim of this thesis was to identify the relationship between R\&D investment and firm-level productivity on Russian data by estimating a hierarchy of production-function models that progress from naive benchmarks to a dynamic endogenous-innovation specification. The empirical design combines a $691{,}644$-observation panel drawn from the RFSD (2011-2023), restricted to the eight production sectors retained after cleaning, with panel-bootstrap inference and a full cross-sectoral replication of every estimator.


At a deliberately generic level, the analysis supports four conclusions. (a)~Firstly, once simultaneity and unobserved heterogeneity are addressed, R\&D investment is associated with a positive and statistically significant productivity gain; the cross-sectional association, by contrast, is dominated by selection and firm fixed effects. (b)~Secondly, the return to R\&D is non-linear and state-dependent: within the observed support of firm productivity, it is the largest for firms operating below the sector productivity median and attenuates towards the upper quantiles, a pattern consistent with the distance-to-frontier prediction of step-by-step Schumpeterian competition and with the catch-up role of innovation in transition economies. This sign of the cross-derivative is the opposite of the predominant complementarity reported by \textcite{doraszelski2013r} for Spanish manufacturing, a contrast that the Discussion attributes to the broader sectoral coverage, the dominance of non-innovators, and the post-sanctions regime in the Russian panel. The exact magnitude and sign at the upper tail are imprecisely identified in the panel bootstrap and are therefore reported as a declining profile rather than a sharply negative one. (c)~Thirdly, the choice of estimator has first-order consequences: moving from reduced-form to structural control-function estimators reattributes a meaningful share of apparent productivity inertia to the cumulative effect of past R\&D, rather than to an exogenous Markov drift. (d)~Fourthly, the cross-sectoral analysis exposes profound structural heterogeneity in production technologies and innovation returns. A notable structural transition is evident in the information and communication sector, as human capital has displaced intermediate materials to become the predominant determinant of value-added. Furthermore, the pervasive near-zero contemporaneous output elasticity of physical capital across all analyzed industries highlights severe rigidities and potential under-utilization in the Russian fixed-asset stock. Most strikingly, the highest marginal returns to R\&D emerge not in mature high-tech manufacturing, but in lagging infrastructure networks such as transportation and utilities. This paradox confirms that basic technological adoption in low-baseline sectors yields outsized catch-up gains, whereas mature industries are already operating in the region of diminishing marginal returns to innovation.


The theoretical contribution is twofold. First, the study operationalizes the endogenous-Markov view of productivity for a setting that spans a full cross-section of production sectors in a large transition economy, extending the external validity of the \textcite{doraszelski2013r} framework beyond its original Spanish benchmark. Second, by estimating the complete degree-3 polynomial, with all ten interaction terms, including the three cross-interactions $\gamma_5$, $\gamma_8$, $\gamma_9$, the analysis provides a direct empirical test of whether the cross-derivative $\partial^2 g/\partial\omega\partial r$ is homogeneous. Notably, the evidence that it is not, the cross-derivative is uniformly negative on the in-sample support and its absolute magnitude grows monotonically across $\omega$-quantiles, so that the first-derivative marginal effect of lagged R\&D decays from $+0.008$ at the bottom decile of within-sample productivity to $+0.001$ at the top decile while remaining positive throughout, is difficult to reconcile with uniform-return models of innovation and supports the distance-to-frontier reading of \textcite{aghion1990model} on the Russian production-economy panel; the sign reversal relative to the complementarity reported by \textcite{doraszelski2013r} for Spanish manufacturing is itself an empirical contribution of the present analysis.


For economic policy, the heterogeneity result translates into a concrete prescription: broad-based R\&D subsidies are dominated by targeted ones, and the marginal productivity payoff of an additional unit of R\&D support is maximised when allocated to firms below the productivity median of their sector. The persistence of the innovator--non-innovator productivity gap during the post-2014 sanctions period is, separately, consistent with domestic innovation partially substituting for restricted access to foreign technology; this observation is descriptive rather than causal but is directly relevant to current debates on import-substitution industrial policy.


Four limitations bound the scope of the findings. First, the R\&D variable is derived from line 1120 of the statutory balance sheet and therefore measures only capitalized R\&D; firms that expense innovation spending directly or report it informally are under-counted, which attenuates the estimated coefficients. Second, the panel does not distinguish process from product innovation and contains no patent or survey-based innovation measures. Third, the \textcite{olley1992dynamics} and \textcite{levinsohn2003estimating} frameworks include a selection stage that can accommodate endogenous firm exit, but this stage was not implemented in the DJ pipeline used here; as a result, the DJ estimates may carry residual survivorship bias in sectors with high exit rates. Fourth, the sample is confined to the eight production sectors of the RFSD, so extrapolation to services or primary-resource extraction is not warranted.


Four directions for further work follow naturally. (i)~Incorporating an exit-selection correction into the DJ estimator, paralleling the \textcite{olley1992dynamics} survival stage, would close the remaining selection gap. (ii)~Re-estimating the degree-3 polynomial on richer innovation proxies, patent records, innovation-survey variables, or direct measures of intangible investment, would relax the line-1120 measurement constraint. (iii)~Modelling spatial knowledge spillovers and inter-firm technology absorption would provide deeper micro-foundations for the heterogeneous cross-derivative documented here. (iv)~Extending the temporal coverage past 2023 would allow a direct test of how the post-2022 sanctions regime reshapes the endogenous component of productivity. Each extension preserves the dynamic structural logic of the present work while relaxing one of its binding assumptions.

% --- REFERENCES ---
\clearpage
\printbibliography[title={References}]

%%
%% APPENDICES
%%
\appendixhead{Data and Code Availability}
\label{app:replication}

All data processing scripts, estimation code, bootstrap routines, and figure/table generation notebooks underlying the empirical analysis reported in this thesis are publicly available in the project's GitHub repository at \href{https://github.com/Ivan-zamyatin}{\faGithub : Ivan-zamyatin}. The repository contains the full replication pipeline, including the RFSD download step, the cleaning and trimming protocol, the six production-function estimators (OLS, FE, FE+time, OP, LP, DJ), the clustered panel bootstrap, and the code used to generate every figure and table in the main text and in this appendix.


\appendixhead{Anomaly Identification and Filtering Protocol}
\label{app:cleaning}

Prior to applying any cleaning filters, an exploratory data analysis (EDA) was conducted on the deflated panel to characterise the distributional properties of the key variables, to identify structural anomalies, and to inform the design of the subsequent filtering protocol.

\begin{figure}[htbp]
   \centering
   \includegraphics[width=0.8\textwidth]{figure_01.png}
   \caption{Share of companies with R\&D investments by sector. The professional, scientific, and technical activities sector shows the highest concentration of innovators.}
   \label{fig:rd_share_by_sector}
\end{figure}

The sectoral composition of innovation activity is examined in Figure~\ref{fig:rd_share_by_sector} \hypertarget{target:rd_share_by_sector}{}, which presents the share of firms reporting positive R\&D expenditure alongside the distribution of revenue across the ten most represented sectors. Professional, scientific and technical activities exhibit the highest concentration of innovating firms, followed by manufacturing, whereas the vast majority of sectors display negligible formal R\&D engagement. The accompanying revenue boxplots reveal substantial cross-sectoral heterogeneity in firm size, with mining and manufacturing firms occupying the upper tail of the revenue distribution. Figure~\ref{fig:log_distributions} \hypertarget{target:log_distributions}{} depicts the $\ln(\text{Value} + 1)$ distributions of the five core financial variables. Revenue, cost of sales, and labour expenditure exhibit approximately symmetric, bell-shaped distributions with small tail near zero in logarithmic space, consistent with a log-normal generating process. Fixed assets display a broader and more dispersed distribution, whilst R\&D expenditure is dominated by a pronounced mass at zero, with a thin right tail representing the small cohort of active innovators. The median of R\&D remains at zero, underscoring the extreme concentration of innovation activity.

The temporal stability of the revenue distribution is assessed in Figure~\ref{fig:revenue_distribution_over_time} \hypertarget{target:revenue_distribution_over_time}{}, which displays violin plots of $\ln(\text{Revenue} + 1)$ for each year from 2012 to 2023. The central tendency and interquartile range remain broadly stable over the observation horizon, with a modest rightward drift in the median consistent with nominal growth and inflationary dynamics. The tails of the distribution, however, exhibit non-trivial variation, particularly during the post-2019 period, which may reflect the transition of data management responsibilities to the FNS. The introduction of a simplified electronic reporting interface in 2019 may have encouraged compliance among non-active firms, thereby inflating the number of 'zero-revenue' companies recorded in the statistics.

Taken together, these five pre-cleaning diagnostics confirm the log-normal character of the standard production inputs, the extreme concentration of R\&D activity, and the pronounced cross-sectoral heterogeneity in firm size and innovation engagement. These features motivate both the logarithmic transformation adopted in the production function specification and the targeted filtering protocol described below. Additional distributional and bivariate diagnostics are provided in the Appendix.\hypertarget{target:revenue_dist_sector}{}\hypertarget{target:dist_pl_revenue}{}\hypertarget{target:dist_fixed_assets}{}\hypertarget{target:dist_cost_sales}{}\hypertarget{target:dist_rd_assets}{}\hypertarget{target:dist_labor_costs}{}\hypertarget{target:box_revenue_rd}{}\hypertarget{target:box_assets_rd}{}\hypertarget{target:box_cost_sales_rd}{}\hypertarget{target:box_labor_rd}{}\hypertarget{target:scatter_rev_assets}{}\hypertarget{target:scatter_rev_cost}{}\hypertarget{target:scatter_rev_labor}{}\hypertarget{target:scatter_assets_labor}{}\hypertarget{target:scatter_assets_cost}{}\hypertarget{target:scatter_cost_labor}{}

Beyond the primary metrics, two diagnostic analyses were integrated to assess revenue concentration and ensure that the findings were not skewed by economic anomalies.  Figure~\ref{fig:lorenz_curve} \hypertarget{target:lorenz_curve}{} presents the Lorenz curve for firm revenue, revealing a pronounced departure from the line of perfect equality: a large share of firms generates only a negligible fraction of total revenue, while only a handful of large enterprises account for the bulk of enterprises. This extreme concentration is a well-documented feature of firm-level microdata, but it also signals the likely presence of inactive or near-zero-revenue entities that do not represent genuine production units. To investigate the sectoral composition of these anomalous observations, firms with revenue at or below 100 thousand RUB were flagged as near-zero-revenue entities. Figure~\ref{fig:low_revenue_high_assets} \hypertarget{target:low_revenue_high_assets}{} identifies the ten sectors with the highest number of such firms that nevertheless report significant fixed assets, highlighting the continued dominance of legacy asset-holding structures, particularly in real estate operations and wholesale trade. Figure~\ref{fig:low_revenue_with_rd} \hypertarget{target:low_revenue_with_rd}{} performs an analogous decomposition for near-zero-revenue firms that report positive R\&D expenditure, revealing that professional, scientific and technical activities and manufacturing dominate this category.It was the insights collected from these diagnostic plots that ultimately dictated the architecture of the a multi-criteria screening procedure, providing the empirical basis for its calibration.

To ensure that the dataset is strictly applicable to a standard production function framework, an integrated filtering procedure was implemented. This vectorised operation isolates and removes economically anomalous observations, inactive entities, and transit structures whose inclusion would contort the empirical estimations. The heuristic criteria, applied jointly, were defined as follows:

\begin{itemize}
    \item \textbf{Non-Operational Entities:} Firms reporting zero or negative values in any core production variable (Revenue $\le 0$, Fixed Assets $\le 0$, Cost of Sales $\le 0$, or Labour Cash Flows $\le 0$) were eliminated, as they violate the fundamental assumptions of production activity.
    \item \textbf{Dormant or Empty Shells:} Observations with negligible fixed assets (Assets $\le 1$ thousand RUB) and minimal revenue (Revenue $\le 100$ thousand RUB) were classified as inactive shell companies and removed.
    \item \textbf{Dead Holdings:} Entities retaining massive capital stocks (Assets $> 10{,}000$ thousand RUB) but exhibiting negligible operational output (Revenue $\le 100$ thousand RUB) and minimal workforce expenditures (Labour $\le 5$ thousand RUB) were excluded. These typically represent legacy asset-holding vehicles rather than actively producing firms.
    \item \textbf{Active Shells:} Firms generating observable revenue (Revenue $> 100$ thousand RUB) but reporting near-zero costs (Cost of Sales $< 1$ thousand RUB) and labour expenses (Labour $< 1$ thousand RUB) were removed, as they indicate a lack of genuine value-added production.
    \item \textbf{Ghost Firms:} Extreme transit structures characterised by massive revenue flows (Revenue $> 5{,}000$ thousand RUB) with a complete absence of fundamental production inputs (Assets $\le 1$ and Labour $\le 1$ thousand RUB) were identified and dropped.
\end{itemize}

In addition, sector-specific filters were calibrated for industries with anomalous capital-labour ratios. Financial and Insurance Activities (OKVED section K) observations with labour costs below 10 thousand RUB, Real Estate Operations (L) observations with fixed assets below 50 thousand RUB, and Professional, Scientific and Technical Activities (M) observations exhibiting high revenue (above 5,000 thousand RUB) but negligible labour costs (below 5 thousand RUB) were flagged as structurally anomalous. Moreover, to ensure the robustness of the results, the dataset was refined by removing firms in the professional, scientific, and technical activities category, as their inclusion could introduce significant specification bias.

To mitigate the outsized influence of extreme statistical outliers while keeping meaningful economic variance, an asymmetric trimming procedure was applied to the remaining dataset. A joint masking protocol was utilised for the four continuous production variables: Revenue, Fixed Assets, Cost of Sales, and Labour Cash Flows. For each of these variables, the top 0.5\% (the 0.995 quantile) and the bottom 1\% (the 0.01 quantile) of observations were systematically truncated. The removal of the bottom 1\% is particularly critical to eliminate anomalous near-zero values that would later give rise to instability logarithmic transformations. Observations failing the threshold for any single variable were excluded, ensuring that the retained sample is jointly clean across all production inputs.

Crucially, Research and Development expenditures were deliberately excluded from this trimming matrix. Because the vast majority of firms in the sample report no formal R\&D activity, applying standard upper-bound truncation would unintentionally screen out the most significant legitimate innovators from the dataset. Preserving the extreme right tail of the R\&D distribution is therefore essential for capturing the true variance of technological investment across the sample.

\begin{figure}[htbp]
   \centering
   \includegraphics[width=0.8\textwidth]{figure_02.png}
   \caption{Revenue distribution by sector on a logarithmic scale. Real estate operations and professional activities exhibit distinct median values and wide ranges of outliers.}
   \label{fig:revenue_dist_sector}
\end{figure}

\begin{figure}[htbp]
   \centering
   \includegraphics[width=0.8\textwidth]{figure_03.png}
   \caption{Natural-log distribution of PL\_revenue ($\ln(\text{Value}+1)$). The data shows a near-normal distribution centred around a median of approximately $\ln(\text{Value}+1)\approx 11$ (i.e.\ a revenue level near $e^{11}\approx 60{,}000$ thousand RUB), with a secondary peak at zero.}
   \label{fig:dist_pl_revenue}
\end{figure}

\begin{figure}[htbp]
   \centering
   \includegraphics[width=0.8\textwidth]{figure_04.png}
   \caption{Distribution of fixed assets (B\_fixed\_assets) on a natural-log scale ($\ln(\text{Value}+1)$). The distribution is bimodal, indicating a large group of firms with very low or zero fixed assets.}
   \label{fig:dist_fixed_assets}
\end{figure}

\begin{figure}[htbp]
   \centering
   \includegraphics[width=0.8\textwidth]{figure_05.png}
   \caption{Distribution of cost of sales (PL\_cost\_of\_sales) using the natural-log transformation $\ln(\text{Value}+1)$. The curve closely follows the distribution observed for total revenue.}
   \label{fig:dist_cost_sales}
\end{figure}

\begin{figure}[htbp]
   \centering
   \includegraphics[width=0.8\textwidth]{figure_06.png}
   \caption{Distribution of R\&D assets (B\_research\_development). The extreme left-skewness confirms that only a tiny fraction of the sampled companies report R\&D assets.}
   \label{fig:dist_rd_assets}
\end{figure}

\begin{figure}[htbp]
   \centering
   \includegraphics[width=0.8\textwidth]{figure_07.png}
   \caption{Natural-log distribution of labour costs (CFo\_labor), plotted as $\ln(\text{Value}+1)$. The median sits near $\ln(\text{Value}+1)\approx 9$ (i.e.\ a labour-cost level near $e^{9}\approx 8{,}000$ thousand RUB), following a typical distribution for operating expenses.}
   \label{fig:dist_labor_costs}
\end{figure}

\begin{figure}[htbp]
   \centering
   \includegraphics[width=0.8\textwidth]{figure_10.png}
   \caption{Spearman correlation matrix of key financial indicators. A very strong correlation (0.97) exists between revenue and cost of sales, while R\&D shows weak correlation with other variables.}
   \label{fig:spearman_correlation}
\end{figure}

\begin{figure}[htbp]
   \centering
   \includegraphics[width=0.8\textwidth]{figure_11.png}
   \caption{Comparison of revenue distribution between firms with and without R\&D. Innovators (Has R\&D) exhibit a significantly higher median revenue level compared to regular firms.}
   \label{fig:box_revenue_rd}
\end{figure}

\begin{figure}[htbp]
   \centering
   \includegraphics[width=0.8\textwidth]{figure_12.png}
   \caption{Box plots of fixed assets (log scale) for innovators versus regular firms. The "Has R\&D" group shows a substantially higher concentration of fixed capital.}
   \label{fig:box_assets_rd}
\end{figure}

\begin{figure}[htbp]
   \centering
   \includegraphics[width=0.8\textwidth]{figure_13.png}
   \caption{Distribution of cost of sales by R\&D status. Similar to revenue, innovators maintain higher operating costs, reflecting a larger scale of operations.}
   \label{fig:box_cost_sales_rd}
\end{figure}

\begin{figure}[htbp]
   \centering
   \includegraphics[width=0.8\textwidth]{figure_14.png}
   \caption{Labor costs comparison between innovators and regular firms. Firms investing in R\&D typically support higher total labor expenditures.}
   \label{fig:box_labor_rd}
\end{figure}

\begin{figure}[htbp]
   \centering
   \includegraphics[width=0.8\textwidth]{figure_15.png}
   \caption{Joint distribution of revenue and fixed assets (log scale). A moderate positive correlation (0.48) is observed, with a dense cluster of firms at lower asset levels.}
   \label{fig:scatter_rev_assets}
\end{figure}

\begin{figure}[htbp]
   \centering
   \includegraphics[width=0.8\textwidth]{figure_16.png}
   \caption{Joint distribution of revenue and cost of sales. The extremely strong correlation (0.97) indicates that these two financial metrics scale almost linearly for most firms.}
   \label{fig:scatter_rev_cost}
\end{figure}

\begin{figure}[htbp]
   \centering
   \includegraphics[width=0.8\textwidth]{figure_17.png}
   \caption{Joint distribution of revenue and labor costs. A strong Spearman correlation of 0.76 suggests that revenue growth is closely tied to investment in human resources.}
   \label{fig:scatter_rev_labor}
\end{figure}

\begin{figure}[htbp]
   \centering
   \includegraphics[width=0.8\textwidth]{figure_18.png}
   \caption{Joint distribution of fixed assets and labor costs. The correlation of 0.53 reflects the varying capital-to-labor intensities across the sampled companies.}
   \label{fig:scatter_assets_labor}
\end{figure}

\begin{figure}[htbp]
   \centering
   \includegraphics[width=0.8\textwidth]{figure_19.png}
   \caption{Joint distribution of fixed assets and cost of sales. The relationship shows a moderate correlation (0.47), consistent with the link between capital and production volume.}
   \label{fig:scatter_assets_cost}
\end{figure}

\begin{figure}[htbp]
   \centering
   \includegraphics[width=0.8\textwidth]{figure_20.png}
   \caption{Joint distribution of cost of sales and labor costs. A strong correlation (0.75) highlights that labor is a primary component of operating expenses for the majority of firms.}
   \label{fig:scatter_cost_labor}
\end{figure}

\begin{figure}[htbp]
   \centering
   \includegraphics[width=0.8\textwidth]{figure_21.png}
   \caption{Comparison of R\&D asset values (B\_research\_development) for firms with and without R\&D. The box plot highlights the significant capital investment gap between the two groups.}
   \label{fig:box_rd_assets_comparison}
\end{figure}

\begin{figure}[htbp]
   \centering
   \includegraphics[width=0.8\textwidth]{figure_24.png}
   \caption{Revenue distribution (PL\_revenue) across different economic sectors. Professional, scientific, and technical activities show a wider range of high-revenue outliers.}
   \label{fig:revenue_by_sector_box}
\end{figure}

\begin{figure}[htbp]
   \centering
   \includegraphics[width=0.8\textwidth]{figure_25.png}
   \caption{Joint distribution and correlation between revenue and R\&D assets. A weak correlation (0.13) suggests that revenue size is not the sole determinant of R\&D investment levels.}
   \label{fig:scatter_rev_rd_assets}
\end{figure}

\begin{figure}[htbp]
   \centering
   \includegraphics[width=0.8\textwidth]{poly_sensitivity_OP.png}
   \caption{Sensitivity analysis for the \textcite{olley1992dynamics} method, showing the stability of labor and capital coefficients across different polynomial orders.}
   \label{fig:poly_sensitivity_op}
\end{figure}

\begin{figure}[htbp]
   \centering
   \includegraphics[width=0.8\textwidth]{residuals_hist_OLS.png}
   \caption{Histogram of residuals from the OLS production function estimation, used to assess the normality of the error term.}
   \label{fig:residuals_hist_ols}
\end{figure}

\begin{figure}[htbp]
   \centering
   \includegraphics[width=0.8\textwidth]{residuals_vs_pred_OLS.png}
   \caption{Residuals versus predicted values for the OLS model, providing a diagnostic check for homoscedasticity and model specification.}
   \label{fig:residuals_vs_pred_ols}
\end{figure}

\appendixhead{Supplementary Model Results and Diagnostics}
\label{app:diagnostics}

This appendix collects the full set of specification tests, sensitivity analyses, and diagnostic plots supporting the headline results reported in Table~\ref{tab:main_results}.

\paragraph{OLS residual diagnostics.}
The OLS baseline achieves $R^2 = 0.9359$; all Variance Inflation Factors remain well below the conventional threshold of $10$, so multicollinearity is not a concern. The Breusch--Pagan LM test ($p<0.001$) and the fan-shaped residual plot (Figure~\ref{fig:residuals_vs_pred_ols}) confirm heteroscedasticity, and the Durbin--Watson statistic of $0.937$ indicates positive serial correlation. Shapiro--Wilk rejects residual normality, which is expected at this sample size and does not threaten consistency; inference therefore relies on HC3 robust standard errors clustered at the firm level for OLS and on the clustered panel bootstrap for the structural models.

\paragraph{TFP persistence and GMM fit across models.}
The FE estimators deliver an implausibly low persistence ($\rho \approx 0.15$), a well-known artefact of the within-transformation amplifying measurement error in short panels. OP returns $\rho = 1.01$, effectively a unit root, raising stationarity concerns. LP and DJ yield economically reasonable values of $0.774$ and $0.759$, respectively; the fact that DJ is marginally lower than LP is consistent with the theory: once R\&D is allowed to influence the transition process, part of what LP attributes to passive persistence is reallocated to the active effect of knowledge accumulation.

The Hansen $J$ statistics reject the overidentifying restrictions at conventional levels for both LP ($J=31.18$, $p<0.001$) and DJ ($J=23.57$, $p<0.001$). This is uninformative on its own: at $n>9\times 10^{5}$ the $J$-test rejects even negligible misspecification. The scaled statistics $J/n = 3.4\times10^{-5}$ (LP) and $2.5\times10^{-5}$ (DJ) indicate that per-observation departures from the moment conditions are economically trivial. DJ shows the smaller $J/n$ and is therefore the preferred structural specification. The mean TFP innovation $\mathbb{E}[\xi_{it}]$ is essentially zero for all three moment-based estimators, confirming that the first-stage inversions are properly centred.

\paragraph{Instrument- and proxy-robustness of LP.}
Table~\ref{tab:lp_lag} reports LP estimates under one-lag and two-lag instrument sets. The capital coefficient shifts modestly ($0.0027 \to 0.0066$), the labour coefficient is essentially unchanged ($0.1463 \to 0.1444$), and the Hansen $J$ falls from $31.18$ to $23.95$ with additional overidentifying restrictions. Neither $J/n$ signals material misspecification.

\begin{figure}[htbp]
   \centering
   \includegraphics[width=0.8\textwidth]{poly_sensitivity_LP.png}
   \caption{Sensitivity of LP coefficient estimates to the degree of the polynomial in the control function.}
   \label{fig:poly_sensitivity_lp}
\end{figure}

Figure~\ref{fig:poly_sensitivity_lp} and Table~\ref{tab:poly_sensitivity} \hypertarget{target:poly_sensitivity}{}\hypertarget{target:poly_sensitivity_op}{} report coefficient stability across polynomial degrees two to four in the LP control function. The labour and materials coefficients vary only modestly, and the capital coefficient remains consistently small, confirming that the nonparametric approximation is not driving the main results.

\begin{table}[htbp]
\centering
\caption{Descriptive Statistics of Financial Variables (Levels, Thousands RUB)}
\label{tab:desc_raw}
\begin{tabular}{lrrrrrrr}
\toprule
Variable & $N$ & Mean & Std & Min & Median & Max & Skew. \\
\midrule
  Revenue & 691,645 & 622,233 & 1,791,714 & 342 & 113,115 & 31,546,636 & 7.24 \\
  Fixed Assets & 691,645 & 2,010 & 8,149 & 0 & 171 & 155,831 & 9.55 \\
  Cost of Sales & 691,645 & 505,424 & 1,438,919 & 169 & 92,944 & 24,175,057 & 7.24 \\
  R\&D Expenditure & 691,645 & 8 & 209 & -249 & 0 & 44,446 & 86.75 \\
  Labor Cost & 691,645 & 76,746 & 206,452 & 58 & 15,493 & 3,016,285 & 6.54 \\
\bottomrule
\end{tabular}
\begin{flushleft}
\footnotesize
\textit{Note:} Sample after basic cleaning (before estimation filters).
All monetary variables deflated to constant prices.
\end{flushleft}
\end{table}

\begin{table}[htbp]
\centering
\caption{OLS Specification Tests}
\label{tab:ols_diagnostics}
\begin{tabular}{lr}
\toprule
Statistic & Value \\
\midrule
  $R^2$               & 0.9359 \\
  Adj.\ $R^2$         & 0.9359 \\
  \midrule
  \multicolumn{2}{l}{\textit{Variance Inflation Factors}} \\
  \quad Labor ($\ln L$) & 2.92 \\
  \quad Materials ($\ln M$) & 2.83 \\
  \quad Capital ($\ln K$) & 1.55 \\
  \quad R\&D Stock ($\ln R$) & 1.03 \\
  \midrule
  Breusch--Pagan LM & 54956.75 \\
  \quad $p$-value    & 0.0000 \\
  Durbin--Watson      & 0.996 \\
  Shapiro--Wilk $W$  & 0.8057 \\
  \quad $p$-value    & 0.0000 \\
\bottomrule
\end{tabular}
\begin{flushleft}
\footnotesize
\textit{Note:} Breusch--Pagan test: $H_0$: homoscedasticity.
Durbin--Watson: values near 2 indicate no autocorrelation.
Shapiro--Wilk: subsample $n=5{,}000$; non-normality expected
at large $N$ but bootstrap inference remains valid.
HC3 heteroscedasticity-robust standard errors used.
\end{flushleft}
\end{table}

\begin{table}[htbp]
\centering
\caption{Marginal Effect of R\&D on Expected Productivity (DJ Model, Degree 3)}
\label{tab:dj_marginal}
\begin{tabular}{rr}
\toprule
$\omega_{t-1}$ & $\partial E[\omega_t] / \partial r_{t-1}$ \\
\midrule
  q10 ($\omega = +0.331$) & +0.0084 \\
  q25 ($\omega = +0.551$) & +0.0083 \\
  q50 ($\omega = +0.752$) & +0.0071 \\
  q75 ($\omega = +0.940$) & +0.0051 \\
  q90 ($\omega = +1.173$) & +0.0014 \\
\bottomrule
\end{tabular}
\begin{flushleft}
\footnotesize
\textit{Note:} Degree-3 marginal effect:
$ME(\omega, r) = \gamma_2 + 2\gamma_4 r + \gamma_5 \omega + 3\gamma_7 r^2 + 2\gamma_9 \omega r + \gamma_8 \omega^2$,
evaluated at $\bar{r} = 3.931$ (mean log R\&D among innovators).
Cross-derivative $\partial^2 g / \partial\omega\partial r = \gamma_5 + 2\gamma_8\omega + 2\gamma_9 r$
varies across firms (unlike degree-2, where it is a constant $\gamma_5$).
\end{flushleft}
\end{table}

\begin{table}[htbp]
\centering
\caption{Cross-Derivative Heterogeneity: $\partial^2 g / \partial\omega\partial r$}
\label{tab:dj_cross_deriv}
\begin{tabular}{lr}
\toprule
Statistic & Value \\
\midrule
Formula & $\gamma_5 + 2\gamma_8\omega + 2\gamma_9 r$ \\
$\gamma_5$ & 0.0080 \\
$\gamma_8$ & -0.0124 \\
$\gamma_9$ & 0.0003 \\
\midrule
Mean & -0.0106 \\
Median & -0.0093 \\
Std.\ dev. & 0.0050 \\
\midrule
\% positive & 0.0\% (107 obs.) \\
\% negative & 100.0\% (528,650 obs.) \\
$N$ & 528,757 \\
\bottomrule
\end{tabular}
\begin{flushleft}
\footnotesize
\textit{Note:} The cross-derivative measures whether R\&D and current productivity
are complements (positive) or substitutes (negative) in the transition function.
Under degree-2, this quantity is a constant ($\gamma_5$); under degree-3,
it varies across firms via the cubic interaction terms $\gamma_8$ and $\gamma_9$.
\end{flushleft}
\end{table}

\begin{table}[htbp]
\centering
\caption{Cross-derivative sign distribution by within-sample quantile of firm productivity $\omega$.  Evaluated on the raw (non-centred) structural $\omega$ that the $\gamma$ coefficients were fitted on, so all bins lie inside the observed support of the data.}
\label{tab:dj_cd_by_quantile}
\begin{tabular}{lrrrrr}
\toprule
$\omega$ quantile bin & Obs. & Mean $\omega$ & Mean CD & \% CD$>$0 & \% CD$<$0 \\
\midrule
  {[0-10]} & 52,876 & +0.587 & -0.0065 & 0.2\% & 99.8\% \\
  {[10-25]} & 79,314 & +0.617 & -0.0073 & 0.0\% & 100.0\% \\
  {[25-50]} & 132,189 & +0.666 & -0.0085 & 0.0\% & 100.0\% \\
  {[50-75]} & 132,189 & +0.750 & -0.0105 & 0.0\% & 100.0\% \\
  {[75-90]} & 79,313 & +0.873 & -0.0134 & 0.0\% & 100.0\% \\
  {[90-100]} & 52,876 & +1.157 & -0.0206 & 0.0\% & 100.0\% \\
\bottomrule
\end{tabular}
\end{table}


\begin{table}[htbp]
\centering
\caption{LP Model: 1-Lag vs 2-Lag Instrument Set Comparison}
\label{tab:lp_lag}
\begin{tabular}{lrr}
\toprule
 & 1-Lag & 2-Lag \\
\midrule
Instruments & $k_{-1}, l_{-1}, k, m_{-1}$ & $+ k_{-2}, l_{-2}, m_{-2}$ \\
Moment conditions & 6 & 10 \\
Parameters & 4 & 4 \\
Overid.\ d.f. & 2 & 6 \\
Observations & 407,953 & 407,953 \\
\midrule
$\beta_k$ (capital) & 0.0138 & 0.0163 \\
$\beta_l$ (labor) & 0.1506 & 0.1224 \\
\midrule
GMM criterion & 0.00002467 & 0.00010210 \\
Hansen $J$ & 10.0645 & 41.6519 \\
$J$ $p$-value & 0.0065 & 0.0000 \\
$J/n$ & 0.00002467 & 0.00010210 \\
\bottomrule
\end{tabular}
\begin{flushleft}
\footnotesize
\textit{Note:} Two-step GMM estimates. Hansen $J$-test: $H_0$: instruments valid.
$J/n < 0.001$ indicates negligible misspecification.
1-Lag instruments: $Z = [k_{t-1}, l_{t-1}, k_t, m_{t-1}]$.
2-Lag adds: $[k_{t-2}, l_{t-2}, m_{t-2}]$.
\end{flushleft}
\end{table}

\begin{table}[htbp]
\centering
\caption{Sensitivity of Coefficients to Polynomial Degree in Control Function}
\label{tab:poly_sensitivity}
\begin{tabular}{lrrrr}
\toprule
Specification & $\beta_l$ & $\beta_m$ & $\beta_k$ & RTS \\
\midrule
\multicolumn{5}{l}{\textit{OP}} \\
  \quad Degree 2 & 0.0691 & 0.8830 & 0.0526 & 1.005 \\
  \quad Degree 3 & 0.0690 & 0.8827 & 0.0493 & 1.001 \\
  \quad Degree 4 & 0.0689 & 0.8825 & 0.0472 & 0.999 \\
\midrule
\multicolumn{5}{l}{\textit{LP}} \\
  \quad Degree 2 & 0.2939 & 0.6941 & 0.0021 & 0.990 \\
  \quad Degree 3 & 0.1307 & 0.8414 & 0.0010 & 0.973 \\
  \quad Degree 4 & 0.1326 & 0.8401 & 0.0010 & 0.974 \\
\bottomrule
\end{tabular}
\begin{flushleft}
\footnotesize
\textit{Note:} Polynomial degrees 2, 3, and 4 for the control function $\phi(\cdot)$.
OP uses investment as proxy; LP uses materials as proxy (ACF identification).
Coefficients estimated on the full sample.
\end{flushleft}
\end{table}

\appendixhead{Cross-Sectoral Estimation by Method}
\label{app:industry_methods}

This appendix collects the per-method industry-specific production function estimates discussed in Section~\ref{sec:cross_sectoral}. For each of the six estimators (Pooled OLS, entity Fixed Effects, two-way Fixed Effects, Olley--Pakes, Levinsohn--Petrin with the ACF correction, and Doraszelski--Jaumandreu), the corresponding industry-by-industry coefficient table is reported alongside the matching coefficient-comparison figure. The synthesis across estimators (returns-to-scale comparison and the four headline conclusions) is presented in the main text and is not duplicated here.



\begin{table}[htbp]
\centering
\caption{Industry-Specific Production Function Estimates (Pooled OLS)}
\label{tab:industry_OLS}
\begin{tabular}{>{\raggedright\arraybackslash}p{4.5cm}rrrrr}
\toprule
Industry & $\beta_l$ & $\beta_m$ & $\beta_k$ & $\beta_r$ & RTS \\
\midrule
  Agriculture, forestry, hunting, fishing and fish farming & 0.0126 & 0.9746 & 0.0268 & 0.0201 & 1.014 \\
  Construction & 0.0787 & 0.8461 & 0.0179 & 0.0091 & 0.943 \\
  Information and communication activities & 0.2363 & 0.7131 & 0.0227 & 0.0118 & 0.972 \\
  Manufacturing & 0.0831 & 0.8933 & 0.0066 & 0.0072 & 0.983 \\
  Mining & 0.0552 & 0.9203 & 0.0197 & 0.0176 & 0.995 \\
  Provision of electric energy, gas and steam; air conditioning & 0.0504 & 0.9537 & 0.0004 & 0.0089 & 1.004 \\
  Transportation and storage & 0.0920 & 0.8068 & 0.0281 & 0.0442 & 0.927 \\
  Water supply; sanitation, waste collection and disposal, pollution control activities & 0.0383 & 0.9648 & -0.0060 & 0.0338 & 0.997 \\
\bottomrule
\end{tabular}
\begin{flushleft}
\footnotesize
\textit{Note:} Pooled OLS estimates by OKVED section. Sectors with fewer than 1{,}000 observations excluded. RTS = $\beta_l + \beta_m + \beta_k$. ``---'' indicates coefficient not estimated by this method.
\end{flushleft}
\end{table}



\begin{figure}[htbp]
   \centering
   \includegraphics[width=0.8\textwidth]{industry_coefficients_OLS.png}
   \caption{Estimated labour, capital, and R\&D coefficients across industries from pooled OLS.}
   \label{fig:industry_coefficients_ols}
\end{figure}



\begin{table}[htbp]
\centering
\caption{Industry-Specific Production Function Estimates (Fixed Effects (Entity))}
\label{tab:industry_FE}
\begin{tabular}{>{\raggedright\arraybackslash}p{4.5cm}rrrrr}
\toprule
Industry & $\beta_l$ & $\beta_m$ & $\beta_k$ & $\beta_r$ & RTS \\
\midrule
  Agriculture, forestry, hunting, fishing and fish farming & 0.1100 & 0.8388 & 0.0417 & -0.0021 & 0.991 \\
  Construction & 0.1206 & 0.8255 & 0.0268 & -0.0019 & 0.973 \\
  Information and communication activities & 0.2506 & 0.5736 & 0.0385 & -0.0007 & 0.863 \\
  Manufacturing & 0.1145 & 0.8574 & 0.0177 & 0.0020 & 0.990 \\
  Mining & 0.0966 & 0.8850 & 0.0245 & 0.0135 & 1.006 \\
  Provision of electric energy, gas and steam; air conditioning & 0.1006 & 0.8565 & 0.0144 & -0.0127 & 0.971 \\
  Transportation and storage & 0.2095 & 0.6562 & 0.0372 & 0.0066 & 0.903 \\
  Water supply; sanitation, waste collection and disposal, pollution control activities & 0.1354 & 0.8199 & 0.0113 & -0.0004 & 0.967 \\
\bottomrule
\end{tabular}
\begin{flushleft}
\footnotesize
\textit{Note:} Fixed Effects (Entity) estimates by OKVED section. Sectors with fewer than 1{,}000 observations excluded. RTS = $\beta_l + \beta_m + \beta_k$. ``---'' indicates coefficient not estimated by this method.
\end{flushleft}
\end{table}



\begin{figure}[htbp]
   \centering
   \includegraphics[width=0.8\textwidth]{industry_coefficients_FE.png}
   \caption{Industry-specific coefficients under entity fixed effects, absorbing time-invariant firm heterogeneity.}
   \label{fig:industry_coefficients_fe}
\end{figure}



\begin{table}[htbp]
\centering
\caption{Industry-Specific Production Function Estimates (Two-Way Fixed Effects (Entity + Time))}
\label{tab:industry_FE_TW}
\begin{tabular}{>{\raggedright\arraybackslash}p{4.5cm}rrrrr}
\toprule
Industry & $\beta_l$ & $\beta_m$ & $\beta_k$ & $\beta_r$ & RTS \\
\midrule
  Agriculture, forestry, hunting, fishing and fish farming & 0.1074 & 0.8361 & 0.0448 & -0.0040 & 0.988 \\
  Construction & 0.1266 & 0.8215 & 0.0232 & 0.0049 & 0.971 \\
  Information and communication activities & 0.2545 & 0.5727 & 0.0368 & 0.0010 & 0.864 \\
  Manufacturing & 0.1103 & 0.8569 & 0.0169 & 0.0014 & 0.984 \\
  Mining & 0.0983 & 0.8856 & 0.0246 & 0.0147 & 1.008 \\
  Provision of electric energy, gas and steam; air conditioning & 0.1076 & 0.8526 & 0.0136 & -0.0099 & 0.974 \\
  Transportation and storage & 0.2122 & 0.6550 & 0.0361 & 0.0087 & 0.903 \\
  Water supply; sanitation, waste collection and disposal, pollution control activities & 0.1410 & 0.8195 & 0.0110 & 0.0020 & 0.972 \\
\bottomrule
\end{tabular}
\begin{flushleft}
\footnotesize
\textit{Note:} Two-Way Fixed Effects (Entity + Time) estimates by OKVED section. Sectors with fewer than 1{,}000 observations excluded. RTS = $\beta_l + \beta_m + \beta_k$. ``---'' indicates coefficient not estimated by this method.
\end{flushleft}
\end{table}



\begin{figure}[htbp]
   \centering
   \includegraphics[width=0.8\textwidth]{industry_coefficients_FE_TW.png}
   \caption{Industry-specific coefficients under two-way fixed effects (entity + year), absorbing both firm heterogeneity and common time shocks.}
   \label{fig:industry_coefficients_fe_tw}
\end{figure}

\begin{figure}[htbp]
   \centering
   \includegraphics[width=0.8\textwidth]{industry_tfp_distribution_OLS.png}
   \caption{Distribution of estimated log-TFP ($\omega_{it}$) by industry from pooled OLS, serving as the reduced-form baseline. Horizontal boxplots ordered by median; whiskers extend to 1.5~IQR.}
   \label{fig:industry_tfp_ols}
\end{figure}

\begin{table}[htbp]
\centering
\caption{Industry-Specific Production Function Estimates (Olley-Pakes (1996))}
\label{tab:industry_OP}
\begin{tabular}{>{\raggedright\arraybackslash}p{4.5cm}rrrrr}
\toprule
Industry & $\beta_l$ & $\beta_m$ & $\beta_k$ & $\beta_r$ & RTS \\
\midrule
  Agriculture, forestry, hunting, fishing and fish farming & -0.0257 & 0.9716 & 0.1105 & 0.0337 & 1.056 \\
  Construction & 0.0652 & 0.8617 & 0.0357 & 0.0243 & 0.963 \\
  Information and communication activities & 0.2381 & 0.6545 & 0.0682 & 0.0177 & 0.961 \\
  Manufacturing & 0.0836 & 0.8747 & 0.0388 & 0.0126 & 0.997 \\
  Mining & 0.0171 & 0.9379 & 0.0456 & 0.0281 & 1.001 \\
  Provision of electric energy, gas and steam; air conditioning & 0.0214 & 0.9715 & 0.0192 & -0.0119 & 1.012 \\
  Transportation and storage & 0.0729 & 0.8081 & 0.0786 & 0.0429 & 0.960 \\
  Water supply; sanitation, waste collection and disposal, pollution control activities & -0.0027 & 0.9853 & 0.0101 & 0.0426 & 0.993 \\
\bottomrule
\end{tabular}
\begin{flushleft}
\footnotesize
\textit{Note:} Olley-Pakes (1996) estimates by OKVED section. Sectors with fewer than 1{,}000 observations excluded. RTS = $\beta_l + \beta_m + \beta_k$. ``---'' indicates coefficient not estimated by this method.
\end{flushleft}
\end{table}



\begin{figure}[htbp]
   \centering
   \includegraphics[width=0.8\textwidth]{industry_coefficients_OP.png}
   \caption{Industry-specific elasticities from the \textcite{olley1992dynamics} estimator with investment as the proxy variable.}
   \label{fig:industry_coefficients_op}
\end{figure}

\begin{figure}[htbp]
   \centering
   \includegraphics[width=0.8\textwidth]{industry_tfp_distribution_OP.png}
   \caption{Distribution of estimated log-TFP ($\omega_{it}$) by industry from the \textcite{olley1992dynamics} estimator (restricted to firm-years with strictly positive investment). Horizontal boxplots ordered by median; whiskers extend to 1.5~IQR.}
   \label{fig:industry_tfp_op}
\end{figure}




\begin{table}[htbp]
\centering
\caption{Industry-Specific Production Function Estimates (Levinsohn-Petrin (2003) with ACF (2015))}
\label{tab:industry_LP}
\begin{tabular}{>{\raggedright\arraybackslash}p{4.5cm}rrrrr}
\toprule
Industry & $\beta_l$ & $\beta_m$ & $\beta_k$ & $\beta_r$ & RTS \\
\midrule
  Agriculture, forestry, hunting, fishing and fish farming & 0.0738 & 0.9500 & 0.0051 & 0.0209 & 1.029 \\
  Construction & 0.0831 & 0.8420 & 0.0123 & 0.0099 & 0.937 \\
  Information and communication activities & 0.2732 & 0.6681 & 0.0132 & 0.0046 & 0.954 \\
  Manufacturing & 0.1124 & 0.8778 & 0.0010 & 0.0030 & 0.991 \\
  Mining & 0.0782 & 0.8263 & 0.0235 & 0.0110 & 0.928 \\
  Provision of electric energy, gas and steam; air conditioning & 0.0829 & 0.8426 & 0.0074 & 0.0165 & 0.933 \\
  Transportation and storage & 0.0918 & 0.7010 & 0.0324 & 0.0412 & 0.825 \\
  Water supply; sanitation, waste collection and disposal, pollution control activities & 0.3791 & 0.5757 & 0.0116 & 0.0219 & 0.966 \\
\bottomrule
\end{tabular}
\begin{flushleft}
\footnotesize
\textit{Note:} Levinsohn-Petrin (2003) with ACF (2015) estimates by OKVED section. Sectors with fewer than 1{,}000 observations excluded. RTS = $\beta_l + \beta_m + \beta_k$. ``---'' indicates coefficient not estimated by this method.
\end{flushleft}
\end{table}



\begin{figure}[htbp]
   \centering
   \includegraphics[width=0.8\textwidth]{industry_coefficients_LP.png}
   \caption{Industry-specific elasticities from the \textcite{levinsohn2003estimating} estimator with the ACF correction.}
   \label{fig:industry_coefficients_lp}
\end{figure}

\begin{figure}[htbp]
   \centering
   \includegraphics[width=0.8\textwidth]{industry_tfp_distribution_LP.png}
   \caption{Distribution of estimated log-TFP ($\omega_{it}$) by industry from the \textcite{levinsohn2003estimating} estimator. Horizontal boxplots ordered by median; whiskers extend to 1.5~IQR, with individual outliers shown.}
   \label{fig:industry_tfp_distribution_lp}
\end{figure}


\begin{table}[htbp]
\centering
\caption{Industry-Specific Production Function Estimates (Doraszelski-Jaumandreu (2013))}
\label{tab:industry_DJ}
\begin{tabular}{>{\raggedright\arraybackslash}p{4.5cm}rrrrr}
\toprule
Industry & $\beta_l$ & $\beta_m$ & $\beta_k$ & $\beta_r$ & RTS \\
\midrule
  Agriculture, forestry, hunting, fishing and fish farming & 0.0714 & 0.9500 & 0.0060 & --- & 1.027 \\
  Construction & 0.0828 & 0.8416 & 0.0125 & --- & 0.937 \\
  Information and communication activities & 0.3078 & 0.6639 & 0.0084 & --- & 0.980 \\
  Manufacturing & 0.1157 & 0.8752 & 0.0010 & --- & 0.992 \\
  Mining & 0.0384 & 0.9500 & 0.0068 & --- & 0.995 \\
  Provision of electric energy, gas and steam; air conditioning & 0.0154 & 0.9500 & 0.0012 & --- & 0.967 \\
  Transportation and storage & 0.0796 & 0.7127 & 0.0323 & --- & 0.825 \\
  Water supply; sanitation, waste collection and disposal, pollution control activities & 0.0314 & 0.9500 & 0.0010 & --- & 0.982 \\
\bottomrule
\end{tabular}
\begin{flushleft}
\footnotesize
\textit{Note:} Doraszelski-Jaumandreu (2013) estimates by OKVED section. Sectors with fewer than 1{,}000 observations excluded. RTS = $\beta_l + \beta_m + \beta_k$. ``---'' indicates coefficient not estimated by this method.
\end{flushleft}
\end{table}


\begin{figure}[htbp]
   \centering
   \includegraphics[width=0.8\textwidth]{industry_coefficients_DJ.png}
   \caption{Industry-specific elasticities from the \textcite{doraszelski2013r} dynamic-innovation estimator (degree-3 polynomial transition).}
   \label{fig:industry_coefficients_dj}
\end{figure}


\begin{figure}[htbp]
   \centering
   \includegraphics[width=0.8\textwidth]{industry_tfp_distribution_DJ.png}
   \caption{Distribution of estimated log-TFP ($\omega_{it}$) by industry from the \textcite{doraszelski2013r} estimator. Horizontal boxplots ordered by median; whiskers extend to 1.5~IQR.}
   \label{fig:industry_tfp_distribution_dj}
\end{figure}



\appendixhead{DJ Polynomial-Order Sensitivity}
\label{app:dj_polynomial_order}

This appendix collects the polynomial-order sensitivity results that justify the choice of the degree-3 transition function reported in the main text. The substantive finding of the appendix is the same as in \textcite{doraszelski2013r}: the degree-2 and degree-3 specifications produce numerically close production-function elasticities and returns to scale, so the case for the richer specification rests on its ability to identify heterogeneity in the cross-derivative $\partial^2 g/\partial\omega\partial r$ rather than on differences in point estimates of input elasticities. Table~\ref{tab:dj_degree_comparison} reports the side-by-side parameter estimates and diagnostics; Table~\ref{tab:marginal_effect_comparison} compares the implied marginal effect of lagged R\&D at five within-sample $\omega$-quantiles; and Figure~\ref{fig:dj_marginal_effect_rd_compare} contrasts the marginal-effect surfaces.

\begin{table}[htbp]
\centering
\caption{DJ Transition Function: Degree-2 vs Degree-3 Comparison}
\label{tab:dj_degree_comparison}
\begin{tabular}{lrr}
\toprule
Parameter & Degree 2 & Degree 3 \\
\midrule
\multicolumn{3}{l}{\textit{Production function coefficients}} \\
  $\beta_l$ (labor)     & 0.0916 & 0.1306 \\
  $\beta_m$ (materials) & 0.8803 & 0.8420 \\
  $\beta_k$ (capital)   & 0.0010 & 0.0010 \\
  RTS                   & 0.9729 & 0.9735 \\
\midrule
\multicolumn{3}{l}{\textit{Transition function parameters}} \\
  $\gamma_0$ (constant) & 0.1578 & 0.0098 \\
  $\gamma_1$ ($\omega_{t-1}$) & 0.8177 & 1.1897 \\
  $\gamma_2$ ($r_{t-1}$) & 0.0237 & -0.0004 \\
  $\gamma_3$ ($\omega_{t-1}^2$) & -0.0605 & -0.2809 \\
  $\gamma_4$ ($r_{t-1}^2$) & -0.0001 & 0.0014 \\
  $\gamma_5$ ($\omega_{t-1} \times r_{t-1}$) & -0.0214 & 0.0080 \\
  $\gamma_6$ ($\omega_{t-1}^3$) & --- & 0.0299 \\
  $\gamma_7$ ($r_{t-1}^3$) & --- & -0.0001 \\
  $\gamma_8$ ($\omega_{t-1}^2 \times r_{t-1}$) & --- & -0.0124 \\
  $\gamma_9$ ($\omega_{t-1} \times r_{t-1}^2$) & --- & 0.0003 \\
\midrule
\multicolumn{3}{l}{\textit{Diagnostics}} \\
  GMM criterion     & 0.000061 & 0.000015 \\
  Hansen $J$        & 32.4399 & 8.1413 \\
  $J/n$             & 0.0001 & 0.0000 \\
\midrule
\multicolumn{3}{l}{\textit{Cross-derivative $\partial^2 g/\partial\omega\partial r$}} \\
  Formula            & $\gamma_5$ (constant) & $\gamma_5 + 2\gamma_8\omega + 2\gamma_9 r$ \\
  \% positive        & 0.0\% & 0.0\% \\
  \% negative        & 100.0\% & 100.0\% \\
\bottomrule
\end{tabular}
\begin{flushleft}
\footnotesize
\textit{Note:} Degree-2 imposes six transition parameters; degree-3 uses the
complete polynomial with ten parameters, as in \textcite{doraszelski2013r}.
Under degree-2, the cross-derivative $\partial^2 g/\partial\omega\partial r = \gamma_5$
is a single constant across all firms; under degree-3, it varies with
$(\omega, r)$ through the higher-order interaction terms.
\end{flushleft}
\end{table}

\begin{table}[htbp]
\centering
\caption{Marginal Effect of R\&D: Degree-2 vs Degree-3}
\label{tab:marginal_effect_comparison}
\begin{tabular}{rrr}
\toprule
Sample quantile ($\omega_{t-1}$) & Degree 2 & Degree 3 \\
\midrule
  q10 ($\omega = +0.331$) & +0.0157 & +0.0084 \\
  q25 ($\omega = +0.551$) & +0.0110 & +0.0083 \\
  q50 ($\omega = +0.752$) & +0.0067 & +0.0071 \\
  q75 ($\omega = +0.940$) & +0.0026 & +0.0051 \\
  q90 ($\omega = +1.173$) & -0.0024 & +0.0014 \\
\bottomrule
\end{tabular}
\begin{flushleft}
\footnotesize
\textit{Note:} Marginal effect $\partial E[\omega_t]/\partial r_{t-1}$ evaluated
at the mean log R\&D among innovators ($\bar{r} = 3.931$).
Degree-2: $\gamma_2 + 2\gamma_4 r + \gamma_5 \omega$.
Degree-3 adds: $+ 3\gamma_7 r^2 + 2\gamma_9 \omega r + \gamma_8 \omega^2$.
\end{flushleft}
\end{table}

\begin{figure}[htbp]
   \centering
   \begin{subfigure}[t]{0.49\textwidth}
      \centering
      \includegraphics[width=\linewidth]{dj_marginal_effect_rd_degree2.png}
      \caption{Degree-2 specification (six $\gamma$-parameters).}
      \label{fig:dj_marginal_effect_rd_deg2}
   \end{subfigure}
   \hfill
   \begin{subfigure}[t]{0.49\textwidth}
      \centering
      \includegraphics[width=\linewidth]{dj_marginal_effect_rd_degree3.png}
      \caption{Degree-3 specification (ten $\gamma$-parameters; preferred).}
      \label{fig:dj_marginal_effect_rd_deg3}
   \end{subfigure}
   \caption{Marginal effect of lagged R\&D on expected productivity, $\partial \mathbb{E}[\omega_{it}]/\partial r_{it-1}$, under the degree-2 (left) and degree-3 (right) \textcite{doraszelski2013r} transition functions. Curves are drawn at the 10th, 50th, and 90th percentiles of the within-sample $\omega$-distribution; the shaded band is the 95\% bootstrap CI at the median $\omega$. Panel~(a) imposes a constant cross-derivative $\partial^2 g/\partial\omega\partial r = \gamma_5$ so the three curves are parallel; panel~(b) admits a cross-derivative that varies across $(\omega, r)$ through the cubic interaction terms $\gamma_7 r^2$, $\gamma_8 \omega^2$ and $\gamma_9 \omega r$, producing the fanned-out pattern that is interpreted in the main text.}
   \label{fig:dj_marginal_effect_rd_compare}
\end{figure}

\end{document}